\chapter{The \texorpdfstring{$x - T$}{x − T} map and \texorpdfstring{$2$}{2}-descent}
\label{ch:descent}

Let $K$ be a field in which $2 \ne 0$, let $f \in K[x]$ be squarefree of odd degree $2g + 1$ with
leading coefficient~$c$, let $L = K[x]/(f)$, and let~$\theta$ be the image of~$x$ in~$L$. The
\emph{descent map}, or $x - T$ map, of Section~5 of M.~Stoll, \emph{Chabauty without the
Mordell--Weil group}, in: Algorithmic and Experimental Methods in Algebra, Geometry, and Number
Theory, Springer, 2017, pp.~623--663 [Sto17], is the homomorphism
\[ \mu \colon J(K) \to L^\times/(L^\times)^2\,, \qquad
   [a, b] \mapsto (-c)^{\deg a}\bigl(a(\theta) - a_1(\theta) f_1(\theta)\bigr)(L^\times)^2\,, \]
where~$a$ is monic and~$a_1$ and~$f_1$ are~$a$ and~$f$ divided by their monic gcd; its kernel
is~$2J(K)$. Over a number field, $\mu$ is the basis of $2$-descent: its image lies in the
$2$-Selmer group, a finite group computable from local data, whose order bounds the rank
of~$J(K)$ when~$J(K)$ is finitely generated. This chapter develops the map, its kernel and the
rank bound, and works out two examples over~$\QQ$.

Section~\ref{sec:hv-classes} introduces the classes $[u, v] \in J(K)$ of Mumford ideals for
arbitrary pairs of polynomials and the identities between them. Section~\ref{sec:ds-map}
defines~$\mu$ without valuation theory, from Mumford representatives prime to~$f$ and an explicit
moving lemma, and proves the formula above, including the correction term~$a_1 f_1$.
Sections~\ref{sec:hv-certificates} and~\ref{sec:hv-cantor} put Proposition~5.1 of [Sto17], which
computes a half of a point of~$2J(K)$, into certificate form: a short list of polynomial
identities and divisibilities between the polynomials that the computation produces implies that
the result is a half, and can be checked by polynomial arithmetic. Section~\ref{sec:ds-kernel}
proves $\ker\mu = 2J(K)$ by the existence half of Proposition~5.1 of [Sto17], which ends in such a
certificate. Section~\ref{sec:ds-basic} shows that the image of~$\mu$ lies in the kernel of the
norm, describes~$J(K)[2]$, and proves the compatibility with base change and the rank formula.
Section~\ref{sec:ds-dedekind}, over the fraction field of a Dedekind domain, bounds the image by
the group $A(\Sigma, 2)$ of classes unramified outside a set~$\Sigma$ of primes, and
Section~\ref{sec:ds-number} defines the $2$-Selmer group and the relaxed Selmer groups over a
number field and proves the rank bound, with the Mordell--Weil theorem and Mattuck's description
of~$J(K_v)$ as cited inputs. Section~\ref{sec:ds-examples} treats two curves of genus~$2$
over~$\QQ$: for the first, $J(\QQ) = J(\QQ)[2]$ and~$C(\QQ)$ is determined without the
Mordell--Weil theorem; for the second, $J(\QQ)$ has rank~$1$.

\emph{Standing hypotheses and conventions.} Throughout the chapter, $K$ is a field in which
$2 \ne 0$, and $f \in K[x]$ is squarefree of odd degree $2g + 1$ with leading coefficient~$c$; as in
Section~5 of [Sto17], $f$ need not be monic. As in Chapter~1, $\A{K} = K[x][y]/(y^2 - f)$ is the
coordinate ring, with the hyperelliptic involution~$\iota$, and $J(K) = \Cl(\A{K})$
(Definitions~\ref{def:hypring} and~\ref{def:hypjac}); the results of Chapters~1 and~4 on Mumford
ideals that are used here need only that~$f$ be squarefree of odd degree. Polynomials are elements
of~$K[x]$; $a \mid b$ means divisibility in~$K[x]$, and $b \equiv b' \bmod a$ means $a \mid b - b'$.
For a nonzero polynomial~$a$, $\operatorname{lc}(a)$ is its leading coefficient. The group~$J(K)$ is
written additively; the groups of square classes are written multiplicatively, so that
$\mu(P + Q) = \mu(P)\mu(Q)$. Sets of primes are denoted by~$\Sigma$, as in Remark~4.2 of [Sto17].

\begin{remark}[Formalization]
The standing hypotheses are a \texttt{variable} block: a field \texttt{K} with
\texttt{[NeZero (2 : K)]} and \texttt{f : K[X]} with \texttt{[Fact (Odd f.natDegree)]} and
\texttt{[Fact (Squarefree f)]}, the instances under which \texttt{HypRing K f} is a Dedekind
domain; the genus is not a variable. The ring~$L$ is \texttt{AdjoinRoot f}, $\theta$ is
\texttt{AdjoinRoot.root f} and~$c$ is \texttt{f.leadingCoeff}. As in Chapter~1, the Jacobian
\texttt{HypJac K f} is Mathlib's multiplicative \texttt{ClassGroup}, so sums, doubles, negatives
and~$0$ in~$J(K)$ become products, squares, inverses and~$1$; $2J(K)$ is the range of
\texttt{powMonoidHom 2}, and $J(K)[2]$ its kernel. The group $L^\times/(L^\times)^2$ is
\texttt{Units.ModPow (AdjoinRoot f) 2}. Lean and Mathlib call a set of primes \texttt{S}. Lean
names without a namespace are in \texttt{ChabautyColeman}. The declarations of this chapter are in
the modules \texttt{HyperellipticCurves.Descent.*} and in \texttt{MumfordValue},
\texttt{MumfordClass} and \texttt{JacobianStructure} of \texttt{HyperellipticCurves.Curve}; the
main theorems of Chapters~5 and~6 do not depend on them.
\end{remark}

\section{Classes of Mumford ideals}
\label{sec:hv-classes}

The Mumford ideal $\mathfrak a_{u,v} = (u, y - v)$ of Definition~\ref{def:mumford-ideal} is
defined for all $u, v \in K[x]$. This section names its class in~$J(K)$ and collects the
identities between such classes that the descent map and the certificates of this chapter
consume. None of them needs a coprimality, degree or monicity condition.

\begin{definition}[Mumford classes]
\label{def:hv-class}
\lean{ChabautyColeman.HypRing.y_sub_algebraMap_ne_zero, ChabautyColeman.mumfordIdeal_ne_bot',
  ChabautyColeman.mumfordIdeal_mem_nonZeroDivisors', ChabautyColeman.mumfordClass,
  ChabautyColeman.mumfordClass_eq_mk0, ChabautyColeman.mumfordClass_eq_one_iff}
\leanok
\uses{def:mumford-ideal, def:hypjac}
For $u, v \in K[x]$, the ideal~$\mathfrak a_{u,v}$ is nonzero, since it contains $y - v \ne 0$.
Its class in~$J(K)$ is denoted by~$[u, v]$; so $[u, v] = 0$ if and only if~$\mathfrak a_{u,v}$ is
principal.
\end{definition}

Every point of~$J(K)$ is~$[u, v]$ for exactly one reduced pair $(u, v)$, that is, one with~$u$
monic, $\deg v < \deg u \le g$ and $u \mid v^2 - f$ (Lemmas~\ref{lem:mumford-reduce}
and~\ref{lem:mumford-exists}). For every pair with~$u$ monic and $u \mid v^2 - f$, $[u, v]$ is the
point denoted in the same way in [Sto17], the class of $(u, v) - (\deg u)\infty$ for the effective
divisor~$(u, v)$ with this Mumford representation. The class~$[u, v]$ is defined for every pair,
not only for pairs with $u \mid v^2 - f$; for instance, $[0, v] = 0$, since
$\mathfrak a_{0,v} = (y - v)$ is principal.

\begin{remark}[Formalization]
\texttt{mumfordClass f u v} is \texttt{ClassGroup.mk0} applied to \texttt{mumfordIdeal f u v} with
the proof \texttt{mumfordIdeal\_mem\_nonZeroDivisors' f u v}, which needs no hypothesis $u \ne 0$.
By proof irrelevance it is the class \texttt{ClassGroup.mk0} of Chapter~1 for any other membership
proof (\texttt{mumfordClass\_eq\_mk0}); the class-level lemmas of Chapter~1 apply after a rewrite
with this lemma, and proofs never unfold the definition. \texttt{mumfordClass\_eq\_one\_iff} is
Mathlib's \texttt{ClassGroup.mk0\_eq\_one\_iff} for \texttt{mumfordClass}.
\texttt{HypRing.y\_sub\_algebraMap\_ne\_zero} holds over every nontrivial commutative ring and needs
nothing about~$f$.
\end{remark}

\begin{lemma}[Basic rules]
\label{lem:hv-class-basic}
\lean{ChabautyColeman.mumfordClass_eq_of_dvd_sub, ChabautyColeman.mumfordClass_one,
  ChabautyColeman.mumfordClass_C_mul, ChabautyColeman.mumfordClass_neg,
  ChabautyColeman.mumfordClass_eq_inv_of_dvd_add, ChabautyColeman.mumfordClass_zero_sq}
\leanok
\uses{def:hv-class}
Let $u, v, v' \in K[x]$.
\begin{enumerate}
  \item If $v \equiv v' \bmod u$, then $[u, v] = [u, v']$. Moreover, $[1, v] = 0$, and
    $[\lambda u, v] = [u, v]$ for $\lambda \in K^\times$.
  \item If $u \mid v^2 - f$, then $[u, -v] = -[u, v]$, and $[u, v'] = -[u, v]$ whenever
    $v' \equiv -v \bmod u$. In particular, $2[u, 0] = 0$ if $u \mid f$.
\end{enumerate}
\end{lemma}

\begin{proof}
\leanok
\uses{lem:mumford-norm, lem:iota-norm}
(1) Already the ideals agree: $y - v = (y - v') - (v - v')$ with $v - v' \in (u)$ gives
$\mathfrak a_{u,v} = \mathfrak a_{u,v'}$; $1 \in \mathfrak a_{1,v}$; and $\lambda u$ and~$u$
generate the same ideal. (2) Here $u \ne 0$, as $v^2 - f \ne 0$: $v^2$ is zero or of even degree,
and~$f$ has odd degree. By Lemma~\ref{lem:mumford-norm},
$\mathfrak a_{u,v}\,\mathfrak a_{u,-v} = \mathfrak a_{u,v}\,\iota(\mathfrak a_{u,v}) = (u)$ is
principal, so $[u, -v] = -[u, v]$; the second claim follows with~(1), and the last one is the case
$v = 0$, where $[u, 0] = [u, -0] = -[u, 0]$.

\emph{Formalization.} The ideal-level facts in~(1) and $v^2 - f \ne 0$ are those of Chapter~1
(\texttt{mumfordIdeal\_eq\_of\_dvd\_sub}, \texttt{mumfordIdeal\_one},
\texttt{mumfordIdeal\_C\_mul}, \texttt{sq\_sub\_ne\_zero}), and the inversion in~(2) is
\texttt{classGroup\_mk0\_mumfordIdeal\_map\_star}. In multiplicative notation, (2) consists of
\texttt{mumfordClass\_neg} (the class of $(u, -v)$ is the inverse of that of $(u, v)$),
\texttt{mumfordClass\_eq\_inv\_of\_dvd\_add} (the same for $(u, v')$ under $u \mid v + v'$) and
\texttt{mumfordClass\_zero\_sq} ($[u, 0]^2 = 1$).
\end{proof}

\begin{lemma}[A B\'ezout identity]
\label{lem:hv-bezout}
\lean{Squarefree.span_eq_top_of_mul_dvd_sq_sub}
\leanok
Let $u_1, u_2, v \in K[x]$ with $u_1 u_2 \mid v^2 - f$. Then~$u_1$, $u_2$ and~$v$ generate the
unit ideal of~$K[x]$: there are $s_1, s_2, s_3 \in K[x]$ with $s_1 u_1 + s_2 u_2 + s_3 v = 1$.
\end{lemma}

\begin{proof}
\leanok
Let~$t$ generate the ideal $(u_1, u_2, v)$ of the principal ideal domain~$K[x]$. Then~$t^2$
divides~$v^2$ and~$u_1 u_2$, hence also $f = v^2 - (v^2 - f)$. Since~$f$ is squarefree, $t$~is a
unit.

\emph{Formalization.} The statement is the equation of ideals $(u_1, u_2, v) = (1)$, for elements
$f, u_1, u_2, v$ of a commutative ring in which every finitely generated ideal is principal
(\texttt{IsBezout}); the proof applies the definition of \texttt{Squarefree f} to a generator,
with no factorization into primes. The triple $(u, 2v, w)$ with $v^2 - f = uw$ in the proof of
Lemma~\ref{lem:mumford-norm} is a special case, as~$2$ is a unit.
\end{proof}

\begin{lemma}[Factorization]
\label{lem:hv-factor}
\lean{ChabautyColeman.mumfordIdeal_mul_mumfordIdeal, ChabautyColeman.mumfordClass_mul,
  ChabautyColeman.mumfordClass_mul_of_dvd_sub, ChabautyColeman.mumfordClass_pow,
  ChabautyColeman.mumfordClass_sq_mul}
\leanok
\uses{def:hv-class}
Let $v \in K[x]$.
\begin{enumerate}
  \item If $u_1, u_2 \in K[x]$ with $u_1 u_2 \mid v^2 - f$, then
    $\mathfrak a_{u_1,v}\,\mathfrak a_{u_2,v} = \mathfrak a_{u_1 u_2, v}$; hence
    $[u_1 u_2, v] = [u_1, v] + [u_2, v]$.
  \item In~(1), $[u_1 u_2, v] = [u_1, v_1] + [u_2, v_2]$ for all $v_1, v_2 \in K[x]$ with
    $v \equiv v_1 \bmod u_1$ and $v \equiv v_2 \bmod u_2$.
  \item If $u \in K[x]$ and $n \ge 0$ with $u^n \mid v^2 - f$, then $[u^n, v] = n[u, v]$. In
    particular, if $A = a_0^2 a$ and $A \mid B^2 - f$, then $[A, B] = 2[a_0, B] + [a, B]$.
\end{enumerate}
\end{lemma}

\begin{proof}
\leanok
\uses{lem:hv-bezout, def:cantor-map, lem:cantor-correct, lem:hv-class-basic}
(1) By Lemma~\ref{lem:hv-bezout}, and since~$2$ is a unit, there are $s_1, s_2, s_3 \in K[x]$
with $s_1 u_1 + s_2 u_2 + s_3 (v + v) = 1$. Cantor's composition in certificate form
(Lemma~\ref{lem:cantor-correct}), with $v_1 = v_2 = v$, gives
$\mathfrak a_{u_1,v}\,\mathfrak a_{u_2,v} = \mathfrak a_{u_1 u_2, V}$, where
\[ V = (s_1 u_1 + s_2 u_2) v + s_3 (v^2 + f) = (1 - 2 s_3 v) v + s_3 (v^2 + f)
     = v - s_3 (v^2 - f)\,. \]
So $V \equiv v \bmod u_1 u_2$, and $\mathfrak a_{u_1 u_2, V} = \mathfrak a_{u_1 u_2, v}$, as in
the proof of Lemma~\ref{lem:hv-class-basic}\,(1). (2) follows from~(1) and
Lemma~\ref{lem:hv-class-basic}\,(1). (3) By induction on~$n$: the case $n = 0$ is
Lemma~\ref{lem:hv-class-basic}\,(1), and for $n \ge 1$ apply~(1) to $u_1 = u$ and
$u_2 = u^{n-1}$. The last claim is~(1) for $u_1 = a_0^2$ and $u_2 = a$, followed by the case
$n = 2$.

\emph{Formalization.} No coprimality of~$u_1$ and~$u_2$ is assumed: a common prime factor~$\pi$
satisfies $\pi^2 \mid v^2 - f$, so $\pi \nmid f$, and the case $u_1 = u_2$ gives
$[u^2, v] = 2[u, v]$. The ideal-level statement takes \texttt{Squarefree f} and $2 \ne 0$ as
explicit hypotheses, the class-level ones the instances. In \texttt{mumfordClass\_pow} the
exponent~$n$ is implicit, as the hypothesis $u^n \mid v^2 - f$ determines it. The last claim
of~(3), \texttt{mumfordClass\_sq\_mul}, is the reduction to squarefree~$a$ in Section~5 of [Sto17]:
$[A, B] \in 2J(K)$ if and only if $[a, B] \in 2J(K)$, and the halves of~$[A, B]$ are the points
$[a_0, B] + Q$ with $2Q = [a, B]$.
\end{proof}

\begin{lemma}[Finite products]
\label{lem:hv-split}
\lean{ChabautyColeman.mumfordClass_prod_of_dvd_sub}
\leanok
\uses{def:hv-class}
Let $n \ge 0$ and $u_1, \dots, u_n, v, v_1, \dots, v_n \in K[x]$ with
$u_1 \cdots u_n \mid v^2 - f$ and $v \equiv v_i \bmod u_i$ for $i = 1, \dots, n$. Then
$[u_1 \cdots u_n, v] = [u_1, v_1] + \dots + [u_n, v_n]$. In particular, if
$b, d_f, d_a, w_1, R, w, r \in K[x]$ satisfy $w = d_f d_a w_1$, $w \mid r^2 - f$,
$r \equiv 0 \bmod d_f$, $r \equiv b \bmod d_a$ and $r \equiv R \bmod w_1$, then
$[w, r] = [d_f, 0] + [d_a, b] + [w_1, R]$.
\end{lemma}

\begin{proof}
\leanok
\uses{lem:hv-factor, lem:hv-class-basic}
By induction on~$n$: the case $n = 0$ is Lemma~\ref{lem:hv-class-basic}\,(1), and for $n \ge 1$,
Lemma~\ref{lem:hv-factor}\,(2), applied to~$u_1 \cdots u_{n-1}$ and~$u_n$, reduces the claim to
the case~$n - 1$. The last claim is the case $n = 3$.

\emph{Formalization.} The factors are indexed by a finite set, and no coprimality is assumed; the
case of distinct linear factors~$u_i$ is \texttt{mk0\_mumfordIdeal\_prod}. The last
claim has no declaration of its own: the proof of Theorem~\ref{thm:hv-general} splits off~$w_1$
and then~$d_f$ and~$d_a$ by Lemma~\ref{lem:hv-factor} (\texttt{mumfordClass\_mul},
\texttt{mumfordClass\_mul\_of\_dvd\_sub}).
\end{proof}

\begin{lemma}[Changing signs on a factor]
\label{lem:hv-twist}
\lean{ChabautyColeman.dvd_sq_sub_of_dvd_sub_of_dvd_add,
  ChabautyColeman.mumfordClass_mul_eq_mul_sq}
\leanok
\uses{def:hv-class}
Let $a_+, a_-, b, \tilde b \in K[x]$ and $a = a_+ a_-$ with $a \mid b^2 - f$,
$\tilde b \equiv b \bmod a_+$ and $\tilde b \equiv -b \bmod a_-$. Then $a \mid \tilde b^2 - f$
and $[a, b] = [a, \tilde b] + 2[a_-, b]$.
\end{lemma}

\begin{proof}
\leanok
\uses{lem:hv-factor, lem:hv-class-basic}
As $\tilde b^2 - f = (b^2 - f) + (\tilde b - b)(\tilde b + b)$ and $a_+ a_-$ divides the product,
$a \mid \tilde b^2 - f$. By Lemma~\ref{lem:hv-factor}, $[a, b] = [a_+, b] + [a_-, b]$ and
$[a, \tilde b] = [a_+, b] + [a_-, -b]$; and $[a_-, -b] = -[a_-, b]$ by
Lemma~\ref{lem:hv-class-basic}\,(2), since $a_- \mid b^2 - f$.

\emph{Formalization.} The lemma does not assume that~$a_+$ and~$a_-$ are coprime. The
divisibility, \texttt{dvd\_sq\_sub\_of\_dvd\_sub\_of\_dvd\_add}, holds in every commutative ring.
\end{proof}

\begin{lemma}[Common factors]
\label{lem:hv-cancel}
\lean{ChabautyColeman.mumfordClass_mul_mumfordClass_eq_of_dvd_add}
\leanok
\uses{def:hv-class}
Let $t, p_1, p_2, v_1, v_2 \in K[x]$ with $t p_1 \mid v_1^2 - f$, $t p_2 \mid v_2^2 - f$ and
$t \mid v_1 + v_2$. Then $[t p_1, v_1] + [t p_2, v_2] = [p_1, v_1] + [p_2, v_2]$.
\end{lemma}

\begin{proof}
\leanok
\uses{lem:hv-factor, lem:hv-class-basic}
By Lemma~\ref{lem:hv-factor}\,(1), $[t p_i, v_i] = [t, v_i] + [p_i, v_i]$ for $i = 1, 2$. Since
$t \mid v_1^2 - f$ and $v_2 \equiv -v_1 \bmod t$, Lemma~\ref{lem:hv-class-basic}\,(2) gives
$[t, v_2] = -[t, v_1]$.

\emph{Formalization.} No coprimality is assumed; the hypothesis $t \mid v_1 + v_2$ gives the
name its suffix \texttt{\_of\_dvd\_add}.
\end{proof}

At the level of ideals, the same argument (Lemmas~\ref{lem:hv-factor}\,(1)
and~\ref{lem:mumford-norm}) gives
\[ \mathfrak a_{tp_1,v_1}\,\mathfrak a_{tp_2,v_2} =
   (t)\,\mathfrak a_{p_1,v_1}\,\mathfrak a_{p_2,v_2}\,, \]
as $\mathfrak a_{t,v_1}\,\mathfrak a_{t,v_2} = \mathfrak a_{t,v_1}\,\iota(\mathfrak a_{t,v_1})
= (t)$: the common factor~$t$ collects the points of the first effective divisor whose images
under~$\iota$ lie on the second, and Cantor's composition removes them together with their images
as the divisor of the function~$t$.

\begin{lemma}[Cantor's reduction steps, certificate form]
\label{lem:hv-reduce}
\lean{ChabautyColeman.mumfordClass_eq_of_sq_sub_eq,
  ChabautyColeman.mumfordClass_eq_of_forall_sq_sub_eq}
\leanok
\uses{def:hv-class}
Let $u, v, u', v' \in K[x]$ and $\lambda \in K$ with $v^2 - f = \lambda u u'$ and
$v' \equiv -v \bmod u'$. Then $[u, v] = [u', v']$. In particular, if $n \ge 0$,
$u_0, v_0, \dots, u_n, v_n \in K[x]$ and $\lambda_0, \dots, \lambda_{n-1} \in K$ satisfy
$v_i^2 - f = \lambda_i u_i u_{i+1}$ and $v_{i+1} \equiv -v_i \bmod u_{i+1}$ for $0 \le i < n$,
then $[u_0, v_0] = [u_n, v_n]$.
\end{lemma}

\begin{proof}
\leanok
\uses{lem:mumford-divisor-y, lem:hv-class-basic}
Since $v^2 - f \ne 0$ (proof of Lemma~\ref{lem:hv-class-basic}\,(2)), $\lambda \ne 0$.
Lemma~\ref{lem:mumford-divisor-y}, applied with $w = \lambda u'$, gives
$[u, v] = [\lambda u', -v]$; by Lemma~\ref{lem:hv-class-basic}\,(1), this is $[u', -v]$, which is
$[u', v']$. The last claim follows by induction on~$n$.

\emph{Formalization.} The step of Lemma~\ref{lem:mumford-divisor-y} used is
\texttt{mk0\_mumfordIdeal\_eq\_mk0\_cofactor\_neg}, with \texttt{mumfordIdeal\_C\_mul};
$\lambda \ne 0$ is not a hypothesis. Unlike the normalized reduction step of
Lemma~\ref{lem:cantor-correct}, which makes the cofactor monic and takes the remainder of~$-v$
modulo it, the lemma allows the cofactor divided by any constant and replaces the remainder by a
congruence. In the second statement the chain consists of sequences $\mathbb N \to K[x]$ and
$\mathbb N \to K$, of which only the values at $0, \dots, n$ (at $0, \dots, n - 1$ for~$\lambda$)
matter, with one hypothesis for all $i < n$ for each condition.
\end{proof}

In Cantor's algorithm, $u'$ is the cofactor $(v^2 - f)/u$ divided by its leading
coefficient~$\lambda$, and~$v'$ is the remainder of~$-v$ modulo~$u'$; if $\deg v < \deg u$ and
$\deg u > g$, then $\deg u' < \deg u$, so that the algorithm terminates
(Lemma~\ref{lem:mumford-reduce}). The lemma needs neither the degree bound nor the normalization.

\section{The descent map}
\label{sec:ds-map}

\begin{setting}
\label{set:ds-setting}
Let $L = K[x]/(f)$, let $\theta \in L$ be the image of~$x$, and let
$L^\square = L^\times/(L^\times)^2$ and $K^\square = K^\times/(K^\times)^2$. Let
$N \colon L^\square \to K^\square$ be the homomorphism induced by the norm~$N_{L/K}$. For a nonzero
$a \in K[x]$, let $d = \gcd(a, f)$ be the monic gcd, and $a = d a_1$, $f = d f_1$, as in
Section~5 of [Sto17].
\end{setting}

Since~$f$ is squarefree, the Chinese remainder theorem gives an isomorphism
$L \cong \prod_p L_p$, $L_p = K[x]/(p)$, over the monic irreducible factors~$p$ of~$f$, which
maps~$\theta$ to the family of the images~$\theta_p$ of~$x$. An element of~$L$ is a unit if and
only if all its components are nonzero, and a unit is a square if and only if all its components
are squares. For $h \in K[x]$, the component of~$h(\theta)$ at~$p$ is~$h(\theta_p)$, which
vanishes if and only if $p \mid h$.

\begin{definition}[The value of a Mumford polynomial]
\label{def:ds-value}
\lean{ChabautyColeman.mumfordValue}
\leanok
For a nonzero $a \in K[x]$ let
\[ m(a) = (-c)^{\deg a}\operatorname{lc}(a)^{-1}\bigl(a(\theta) - a_1(\theta) f_1(\theta)\bigr)
   \in L\,. \]
If~$a$ is prime to~$f$, then $d = 1$, $f_1 = f$, and
$m(a) = (-c)^{\deg a}\operatorname{lc}(a)^{-1} a(\theta)$.
\end{definition}

For monic~$a$, the class of~$m(a)$ is the class $\tilde\mu(a)$ of Section~5 of [Sto17]. The factor
$\operatorname{lc}(a)^{-1}$ makes $m(\lambda a) = m(a)$ for $\lambda \in K^\times$ (multiplying~$a$
by~$\lambda$ multiplies~$a_1$ by~$\lambda$ and leaves~$d$ and~$f_1$ unchanged), in accordance with
$[\lambda a, b] = [a, b]$ (Lemma~\ref{lem:hv-class-basic}\,(1)). The value depends only on~$a$; a
second polynomial~$b$ enters only through the hypothesis $a \mid b^2 - f$ of the next lemma, which
makes~$m(a)$ a unit.

\begin{lemma}[Units and multiplicativity]
\label{lem:ds-value}
\lean{ChabautyColeman.isUnit_mumfordValue, ChabautyColeman.mumfordValue_mul,
  Squarefree.dvd_of_dvd_of_dvd_sq_sub, Squarefree.isRelPrime_of_dvd_of_mul_dvd_sq_sub}
\leanok
\uses{def:ds-value}
\begin{enumerate}
  \item Let $a, b \in K[x]$ with $a \mid b^2 - f$. Then~$d$ is squarefree, $d$ and~$a_1$ are
    coprime, $d \mid b$, and~$m(a)$ is a unit of~$L$: at a monic irreducible factor~$p$ of~$f$,
    its component is $-(-c)^{\deg a}\operatorname{lc}(a)^{-1}a_1(\theta_p)f_1(\theta_p)$ if
    $p \mid d$, and $(-c)^{\deg a}\operatorname{lc}(a)^{-1}a(\theta_p)$ if $p \mid f_1$; both are
    nonzero.
  \item If $a, a' \in K[x]$ are nonzero and prime to~$f$, then $m(aa') = m(a)m(a')$.
\end{enumerate}
\end{lemma}

\begin{proof}
\leanok
(1) $a \ne 0$, as $b^2 - f \ne 0$ ($f$ has odd degree), and~$d$ divides the squarefree~$f$. Every
prime factor of~$d$ divides~$f$ and $b^2 - f$, hence~$b$; so the squarefree~$d$ divides~$b$. If a
prime~$p$ divided both~$d$ and~$a_1$, then $p^2 \mid a \mid b^2 - f$ and $p \mid b$, so
$p^2 \mid f$, which is impossible. Let $p \mid f$ be monic irreducible. If $p \mid d$, then
$p \mid a$, $p \nmid a_1$ and $p \nmid f_1$, as~$f$ is squarefree. If $p \mid f_1$, then
$p \nmid d$, so $p \nmid a$, as a common factor of~$a$ and~$f$ divides~$d$; and
$f_1(\theta_p) = 0$. (2) $aa'$ is prime to~$f$, and degrees and leading coefficients are
multiplicative.

\emph{Formalization.} \texttt{mumfordValue f a} is the image in \texttt{AdjoinRoot f} of
$(-c)^{\deg a}\operatorname{lc}(a)^{-1}(a - a_1 f_1)$, with the normalized (monic) gcd
of~$K[x]$, taken classically; its value at $a = 0$ is~$0$, as $\operatorname{lc}(0)^{-1} = 0$ in
Lean. Proofs rewrite with \texttt{mumfordValue\_eq\_of\_mul\_eq}, which computes~$m(a)$ from any
factorizations $a = d a_1$, $f = d f_1$ with monic~$d$ and $a_1$ prime to~$f_1$. The proof of~(1)
needs no decomposition of~$L$: it writes $a - a_1 f_1 = a_1(d - f_1)$, where~$a_1$ is prime to
$f = d f_1$ and $d - f_1$ is prime to $d f_1$
(\texttt{Squarefree.isCoprime\_sub\_mul\_of\_dvd\_sq\_sub}), and an element prime to~$f$ has a unit
image in \texttt{AdjoinRoot f} (\texttt{AdjoinRoot.isUnit\_mk\_iff}). The claims that~$d$
divides~$b$ and is prime to~$a_1$ hold in every commutative ring that is a decomposition monoid.
\end{proof}

\begin{lemma}[Well-definedness]
\label{lem:ds-coprime}
\lean{ChabautyColeman.isSquare_mumfordValue_mul_of_mumfordClass_eq}
\leanok
\uses{def:ds-value, def:hv-class}
Let $a, b, a', b' \in K[x]$ with $a \mid b^2 - f$ and $a' \mid b'^2 - f$, with~$a$ and~$a'$ prime
to~$f$, and with $[a, b] = [a', b']$. Then $m(a)m(a')$ is a square in~$L^\times$.
\end{lemma}

\begin{proof}
\leanok
\uses{lem:mumford-norm, lem:iota-norm, lem:normform, lem:hv-class-basic, lem:ds-value}
By Lemma~\ref{lem:hv-class-basic}\,(2), $[a', -b'] = -[a', b'] = -[a, b]$, so the ideal
$\mathfrak a_{a,b}\,\mathfrak a_{a',-b'}$ is principal, generated by $h = p + qy$ with
$p, q \in K[x]$. By Lemma~\ref{lem:mumford-norm},
$\mathfrak a_{a,b}\,\iota(\mathfrak a_{a,b}) = (a)$ and
$\mathfrak a_{a',-b'}\,\iota(\mathfrak a_{a',-b'}) = (a')$, so the principal ideal of~$\A{K}$
generated by the norm $h\,\iota(h) = p^2 - fq^2 \in K[x]$ (Lemma~\ref{lem:normform}) is generated
by~$aa'$. As~$\A{K}$ is a free $K[x]$-module with basis $1, y$, an ideal of~$K[x]$ is the
contraction of the ideal of~$\A{K}$ that it generates; hence $p^2 - fq^2 = \lambda aa'$ with
$\lambda \in K^\times$. Evaluating at~$\theta$, where~$f$ vanishes, gives
$p(\theta)^2 = \lambda a(\theta)a'(\theta)$, so
\[ m(a)m(a') = (-c)^D\bigl(\lambda\operatorname{lc}(a)\operatorname{lc}(a')\bigr)^{-1}
   p(\theta)^2\,, \qquad D = \deg a + \deg a'\,. \]
The degree of~$p^2$ is even, and that of~$fq^2$ is odd if $q \ne 0$, so the leading term of
$p^2 - fq^2$ is that of~$p^2$ if~$D$ is even and that of~$-fq^2$ if~$D$ is odd. In the first case
$\lambda\operatorname{lc}(a)\operatorname{lc}(a') = \operatorname{lc}(p)^2$, and $(-c)^D$ is a
square; in the second $\lambda\operatorname{lc}(a)\operatorname{lc}(a') =
-c\operatorname{lc}(q)^2$, and
$(-c)^D/(-c\operatorname{lc}(q)^2) = (-c)^{D-1}\operatorname{lc}(q)^{-2}$ is a square. In both
cases $m(a)m(a')$ is a square constant times~$p(\theta)^2$, and~$p(\theta)$ is a unit, as
$a(\theta)$ and~$a'(\theta)$ are (Lemma~\ref{lem:ds-value}\,(1)).

\emph{Formalization.} The norm identity of Lemma~\ref{lem:mumford-norm} carries the
well-definedness; no valuation theory of~$\A{K}$ is needed. The contraction is Mathlib's
\texttt{Ideal.map\_injective\_of\_faithfullyFlat}, as~$\A{K}$ is free over~$K[x]$
(\texttt{associated\_norm\_of\_mumfordIdeal\_mul\_mumfordIdeal\_eq\_span}). Associated polynomials
have the same value (\texttt{mumfordValue\_eq\_of\_associated}), and the case distinction on the
leading term is
\texttt{Polynomial.isSquare\_neg\_leadingCoeff\_pow\_mul\_leadingCoeff\_sq\_sub\_mul\_sq}; the step
that~$p(\theta)$ is a unit is not needed. The conclusion is that $m(a)m(a')$ is a square in
\texttt{AdjoinRoot f}; in~$L^\square$, \texttt{Units.ModPow.mk\_eq\_one\_iff\_isSquare} translates
it.
\end{proof}

\begin{lemma}[Moving, {[Sto17, Section~5]}]
\label{lem:ds-move}
\lean{ChabautyColeman.mumfordClass_eq_of_mul_eq_sub_sq,
  ChabautyColeman.exists_isCoprime_mumfordClass_eq}
\leanok
\uses{def:hv-class}
Let $a, b \in K[x]$ with $a \mid b^2 - f$.
\begin{enumerate}
  \item Let $h, a'' \in K[x]$ with $aa'' = f - (b + ha)^2$. Then $a'' \mid (b + ha)^2 - f$ and
    $[a, b] = [a'', -(b + ha)]$.
  \item For every nonzero $g \in K[x]$ there is $h \in K[x]$ such that the polynomial~$a''$
    of~(1) is prime to~$fg$.
\end{enumerate}
\end{lemma}

\begin{proof}
\leanok
\uses{lem:hv-class-basic, lem:hv-reduce}
(1) The divisibility is clear. By Lemma~\ref{lem:hv-class-basic}\,(1), $[a, b] = [a, b + ha]$;
and since $(b + ha)^2 - f = (-1) \cdot a \cdot a''$, Lemma~\ref{lem:hv-reduce} with
$\lambda = -1$ and $v' = -(b + ha)$ gives $[a, b + ha] = [a'', -(b + ha)]$.

(2) Write $f - b^2 = a a_0''$; then $a'' = a_0'' - 2hb - h^2 a$. For a monic irreducible
factor~$p$ of~$fg$, let~$\theta_p$ be the image of~$x$ in the field $k_p = K[x]/(p)$ and
\[ Q_p(t) = a_0''(\theta_p) - 2b(\theta_p)t - a(\theta_p)t^2 \in k_p[t]\,, \]
so that $a''(\theta_p) = Q_p(h(\theta_p))$. The polynomial~$Q_p$ is not zero: otherwise~$p$
divides~$a$, $b$ and~$a_0''$, hence $p^2 \mid aa_0'' + b^2 = f$, which is impossible. So~$Q_p$ has
at most two roots in~$k_p$, while~$k_p$ has at least three elements, as $2 \ne 0$ in~$K$. Choose
$t_p \in k_p$ with $Q_p(t_p) \ne 0$ for each of the finitely many~$p$, and~$h$ with
$h(\theta_p) = t_p$ for all~$p$, by the Chinese remainder theorem. Then no~$p$ divides~$a''$,
so~$a''$ is prime to~$fg$.

\emph{Formalization.} In~(1), \texttt{k} stands for~$h$. The proof of~(2) avoids the residue
fields: \texttt{Squarefree.exists\_mul\_eq\_sub\_sq\_and\_isCoprime} holds in every principal
ideal domain in which~$2$ is a unit, with any nonzero~$N$ in place of~$fg$. No prime~$p$ divides
the values of $a_0'' - 2bt - at^2$ at all three of $t = 0, 1, -1$, as it would then divide~$a$,
$b$ and~$a_0''$; an induction over the prime factors of~$N$, with the Chinese remainder theorem for
two coprime elements at each step (\texttt{IsCoprime.exists\_dvd\_sub\_and\_dvd\_sub}), gives~$h$.
\end{proof}

In the notation of [Sto17], after Proposition~5.1, $f - b^2 = ac'$ and
$[c' - 2hb - h^2 a, -b - ha] = [a, b]$; so $a'' = c' - 2hb - h^2 a$, and [Sto17] takes~$h$
constant to make~$a''$ prime to~$f$. If~$K$ is infinite, all but finitely many constants
$h = t \in K$ make~$a''$ prime to~$fg$, as in the proof of Lemma~\ref{lem:ds-move}\,(2):
$a''(\theta_p) = Q_p(t)$, and each~$Q_p$ has at most two roots in $K \subseteq k_p$. This remark
is not formalized.

\begin{definition}[The $x - T$ map]
\label{def:ds-mu}
\lean{ChabautyColeman.HypJac.descentMap,
  ChabautyColeman.HypJac.descentMap_mumfordClass_of_isCoprime}
\leanok
\uses{def:ds-value, def:hv-class}
There is a unique homomorphism $\mu \colon J(K) \to L^\square$ with
$\mu([a, b]) = m(a)(L^\times)^2$ for all $a, b \in K[x]$ with $a \mid b^2 - f$ and~$a$ prime
to~$f$.
\end{definition}

\begin{proof}
\leanok
\uses{lem:ds-move, lem:ds-coprime, lem:ds-value, lem:mumford-reduce, lem:hv-factor}
Every point of~$J(K)$ is $[a, b]$ with $a \mid b^2 - f$ (Lemma~\ref{lem:mumford-reduce}), and by
Lemma~\ref{lem:ds-move}\,(2) with $g = 1$ also with~$a$ prime to~$f$. So the formula determines a
map, which does not depend on the representative by Lemma~\ref{lem:ds-coprime}, as~$L^\square$
has exponent~$2$. Let $P = [a, b]$ and $Q = [a', b']$ with~$a$ and~$a'$ prime to~$f$; by
Lemma~\ref{lem:ds-move}\,(2) with $g = a$, the polynomial~$a'$ can be chosen prime to~$a$ as well.
By the Chinese remainder theorem there is~$B$ with $B \equiv b \bmod a$ and
$B \equiv b' \bmod a'$; then $aa' \mid B^2 - f$, as~$a$ and~$a'$ are coprime and each divides
$B^2 - f$, and $[aa', B] = P + Q$ by Lemma~\ref{lem:hv-factor}\,(2). As~$aa'$ is prime to~$f$,
$\mu(P + Q)$ is the class of $m(aa') = m(a)m(a')$ (Lemma~\ref{lem:ds-value}\,(2)).

\emph{Formalization.} \texttt{HypJac.descentMap f} chooses, for each point, a representative prime
to~$f$ (\texttt{HypJac.exists\_isCoprime\_mumfordClass\_eq}, from
\texttt{exists\_mk0\_mumfordIdeal\_reduced} and the moving lemma), and
\texttt{HypJac.descentMap\_mumfordClass\_of\_isCoprime} characterizes it: the value is the class of
any unit~$u$ with $u = m(a)$. Later proofs use only this lemma and
Theorem~\ref{thm:ds-formula}, never the choice. Two units have the same class if and only if their
product is a square (\texttt{Units.ModPow.mk\_eq\_mk\_iff\_isSquare\_mul}), which turns
Lemma~\ref{lem:ds-coprime} into the independence of the choice.
\end{proof}

\begin{theorem}[The formula for~$\mu$, {[Sto17, Section~5]}]
\label{thm:ds-formula}
\lean{ChabautyColeman.HypJac.descentMap_mumfordClass,
  ChabautyColeman.HypJac.descentMap_mumfordClass_X_sub_C}
\leanok
\uses{def:ds-mu, def:ds-value}
For all $a, b \in K[x]$ with $a \mid b^2 - f$, $\mu([a, b])$ is the class of~$m(a)$.
\end{theorem}

\begin{proof}
\leanok
\uses{lem:ds-move, lem:ds-value}
By Lemma~\ref{lem:ds-move}, there are~$h$ and~$a''$ with $aa'' = f - v^2$, $v = b + ha$, such
that~$a''$ is prime to~$f$ and $[a, b] = [a'', -v]$; so~$\mu([a, b])$ is the class of~$m(a'')$.
By Lemma~\ref{lem:ds-value}\,(1), $d \mid b$, so $v = dv_1$ with $v_1 = b/d + ha_1$, and
$a_1a'' = f_1 - dv_1^2$ after cancelling~$d$. Hence the polynomial identity
\[ (a - a_1f_1)\,a'' = f\,(1 + v_1)^2 - (v + f_1)^2\,: \]
indeed $(a - a_1f_1)a'' = aa'' - f_1 \cdot a_1a'' = f - v^2 - f_1^2 + fv_1^2$ and
$(v + f_1)^2 = v^2 + 2fv_1 + f_1^2$. Evaluating at~$\theta$ gives
$\bigl(a(\theta) - a_1(\theta)f_1(\theta)\bigr)a''(\theta) = -s^2$ with $s = (v + f_1)(\theta)$,
so
\[ m(a)m(a'') = -(-c)^D\bigl(\operatorname{lc}(a)\operatorname{lc}(a'')\bigr)^{-1}s^2\,, \qquad
   D = \deg a + \deg a''\,. \]
If $2\deg v > 2g + 1$, then~$D$ is even and
$\operatorname{lc}(a)\operatorname{lc}(a'') = -\operatorname{lc}(v)^2$, so
$m(a)m(a'') = c^D\operatorname{lc}(v)^{-2}s^2$; otherwise $D = 2g + 1$ and
$\operatorname{lc}(a)\operatorname{lc}(a'') = c$, so $m(a)m(a'') = c^{2g}s^2$. In both cases
$m(a)m(a'')$ is a square in~$L^\times$ (both factors are units by
Lemma~\ref{lem:ds-value}\,(1), so~$s$ is one), and the classes of~$m(a)$ and~$m(a'')$ agree.

\emph{Formalization.} The comparison of~$m(a)$ and $m(a'')$ is
\texttt{isSquare\_mumfordValue\_mul\_of\_mul\_eq\_sub\_sq}; the identity is checked by
\texttt{linear\_combination} after substituting $a = da_1$, $f = df_1$ and $v = dv_1$. Only one
move is needed, the moving lemma with~$f$ as the polynomial to avoid.
\texttt{HypJac.descentMap\_mumfordClass\_X\_sub\_C} is the case $a = x - x_0$ for a point
$(x_0, y_0)$ with $f(x_0) \ne 0$, where $m(a) = c(x_0 - \theta)$.
\end{proof}

Section~5 of [Sto17] states this formula and refers to Schaefer for it; the identity above proves
it. For~$a$ prime to~$f$, Theorem~\ref{thm:ds-formula} is the definition.

\begin{corollary}[Values at the points of order~$2$]
\label{cor:ds-torsion-values}
\lean{ChabautyColeman.HypJac.descentMap_mumfordClass_zero_of_mul_eq,
  ChabautyColeman.HypJac.map_liftAlgHom_descentMap_mumfordClass_X_sub_C_zero_of_ne,
  ChabautyColeman.HypJac.map_liftAlgHom_descentMap_mumfordClass_X_sub_C_zero_self}
\leanok
\uses{def:ds-mu}
Let $h \in K[x]$ be a nonzero divisor of~$f$. Then $\mu([h, 0])$ is the class of
\[ (-c)^{\deg h}\operatorname{lc}(h)^{-1}\bigl(h(\theta) - (f/h)(\theta)\bigr)\,, \]
whose component at a monic irreducible factor~$p$ of~$f$ is
$-(-c)^{\deg h}\operatorname{lc}(h)^{-1}(f/h)(\theta_p)$ if $p \mid h$ and
$(-c)^{\deg h}\operatorname{lc}(h)^{-1}h(\theta_p)$ if $p \nmid h$.
\end{corollary}

\begin{proof}
\leanok
\uses{thm:ds-formula}
$h \mid 0^2 - f$. Here $d = \operatorname{lc}(h)^{-1}h$, $a_1 = \operatorname{lc}(h)$ and
$f_1 = \operatorname{lc}(h)f/h$, so $a_1f_1 = \operatorname{lc}(h)^2f/h$. Hence~$m(h)$ differs from
the displayed element by the square factor~$\operatorname{lc}(h)^2$ in the components at the
factors of~$h$ and agrees with it in the others; apply Theorem~\ref{thm:ds-formula}.

\emph{Formalization.} The Lean statement takes the cofactor~$e$ with $he = f$ in place of~$f/h$,
so that no division appears. For a root~$x_0$ of~$f$ in~$K$ and $h = x - x_0$, the components of
the value at the roots~$x_1$ of~$f$ are $\operatorname{lc}(f)(x_0 - x_1)$ for $x_1 \ne x_0$ and
$\operatorname{lc}(f)f'(x_0)$ for $x_1 = x_0$ (the two declarations with
\texttt{map\_liftAlgHom}, for the evaluation $L \to K$ at~$x_1$).
\end{proof}

For monic~$h$ the two elements are equal. The values $\mu([h, 0])$ describe~$\mu$ on~$J(K)[2]$
(Theorem~\ref{thm:ds-two-torsion}).

\begin{remark}[The elliptic case]
\label{rem:ds-ec}
For $g = 1$ and monic~$f$, $\mu$ is the classical $x - T$ map of the elliptic curve
$y^2 = f(x)$. A point $P = (x_0, y_0)$ corresponds to $[x - x_0, y_0]$ (the Abel--Jacobi map of
Definition~\ref{def:hypjac}), and $m(x - x_0) = x_0 - \theta$ if $f(x_0) \ne 0$, while for
$f(x_0) = 0$, where $d = x - x_0$, $a_1 = 1$ and $f_1 = f/(x - x_0)$,
$m(x - x_0) = x_0 - \theta + f_1(\theta)$, whose component at $x - x_0$ is $f'(x_0)$.
\end{remark}

\section{Halving certificates}
\label{sec:hv-certificates}

For a point $P \in 2J(K)$, Proposition~5.1 of [Sto17] computes a point~$Q$ with $2Q = P$ from the
solution of a system of linear congruences. This section and the next turn that computation into
a certificate: Theorems~\ref{thm:hv-generic}, \ref{thm:hv-general} and~\ref{thm:hv-add} and
Corollary~\ref{cor:hv-record} say that a short list of polynomial identities and divisibilities
between the polynomials that the computation produces implies $2Q = P$. These conditions can be
checked by polynomial arithmetic, so neither the linear system, nor the square root~$s$ of the
proposition, nor the minimality of its solution, nor Cantor's algorithm has to be trusted. The
one new idea is Lemma~\ref{lem:hv-principal}: the function~$uy - v$ in the proof of
Proposition~5.1 of [Sto17] generates a Mumford ideal. This turns the computation of the divisor
of~$uy - v$ there into an identity of ideals, so that the halving itself needs no reduction step;
Cantor's composition and reduction enter only through the sum that forms the half when the first
polynomial of~$P$ is not squarefree (Section~\ref{sec:hv-cantor}). The proof of
$\ker\mu = 2J(K)$ in Section~\ref{sec:ds-kernel} ends in a certificate of this section.

\begin{lemma}[The function $uy - v$]
\label{lem:hv-principal}
\lean{ChabautyColeman.dvd_mul_sq_sub_of_dvd_sq_sub_sq_mul,
  ChabautyColeman.span_singleton_mul_y_sub_eq_mumfordIdeal,
  ChabautyColeman.mumfordClass_eq_one_of_sq_sub_eq}
\leanok
\uses{def:mumford-ideal, def:hv-class}
Let $u, v, e \in K[x]$, and let $N = v^2 - u^2 f$ be the norm of~$uy - v$.
\begin{enumerate}
  \item If $eu \equiv 1 \bmod N$, then the principal ideal $(uy - v)$ of~$\A{K}$ is the Mumford
    ideal~$\mathfrak a_{N, ev}$.
  \item Let $A \in K[x]$ and $\lambda \in K$ with $N = \lambda A$ and $eu \equiv 1 \bmod A$. Then
    $A \mid (ev)^2 - f$. If moreover $\lambda \ne 0$, then $[A, ev] = 0$.
\end{enumerate}
\end{lemma}

\begin{proof}
\leanok
\uses{lem:hv-class-basic, lem:mumford-exists}
(1) Write $eu = 1 + tN$ with $t \in K[x]$. Then $uy - v = u (y - ev) + tvN$ lies
in~$\mathfrak a_{N, ev}$. Conversely, $(uy - v)(uy + v) = u^2 f - v^2 = -N$, so~$N$ lies
in~$(uy - v)$, and so does $y - ev = (uy - v)\bigl(e + ty(uy + v)\bigr)$. Hence
$(uy - v) = \mathfrak a_{N, ev}$.
(2) Write $eu = 1 + tA$ with $t \in K[x]$. Then
\[ (ev)^2 - f = e^2 (v^2 - u^2 f) + (e^2 u^2 - 1) f
               = A\bigl(\lambda e^2 + t (eu + 1) f\bigr)\,. \]
If $\lambda \ne 0$, then $eu \equiv 1 \bmod N$, so $[N, ev] = 0$ by~(1), and
$[A, ev] = [\lambda A, ev] = [N, ev]$ by Lemma~\ref{lem:hv-class-basic}\,(1).

\emph{Formalization.} The ideal statement~(1) is stated for the norm~$N$ itself,
without~$\lambda$, and holds over every field, with no hypothesis on~$f$; its proof checks
membership in the Mumford ideal with \texttt{mem\_mumfordIdeal\_iff}. The statements of~(2), which
the certificates use, keep~$\lambda$, called \texttt{lam} as~$\lambda$ is a keyword of Lean; the
class-level one also assumes $\lambda \ne 0$ and passes from~$\lambda A$ to~$A$ by
\texttt{mumfordClass\_C\_mul}. For $u = e = 1$, (1) says that~$(y - v)$ is the Mumford ideal
of~$v^2 - f$; with Lemma~\ref{lem:hv-factor} this gives the ideal identity of
Lemma~\ref{lem:mumford-divisor-y}.
\end{proof}

The second polynomial of the Mumford ideal is $ev \equiv v/u \bmod N$, as it must be: at a
root~$\xi$ of~$N$, the function~$uy - v$ vanishes at the point $(\xi, v(\xi)/u(\xi))$.

\begin{definition}[Halving certificates]
\label{def:hv-cert}
\lean{ChabautyColeman.HalvingCertificate}
\leanok
Let $a, b, w, r \in K[x]$. A \emph{halving certificate} for $(a, b)$ with half $(w, r)$ consists
of polynomials $u, v, e, e' \in K[x]$ and a constant $\lambda \in K^\times$ such that
\begin{description}
  \item[\textmd{(H1)}] $u^2 f = v^2 - \lambda a w^2$,
  \item[\textmd{(H2)}] $v \equiv ub \bmod a$,
  \item[\textmd{(H3)}] $ur \equiv -v \bmod w$, and
  \item[\textmd{(H4)}] $eu + e' a w^2 = 1$.
\end{description}
\end{definition}

In Proposition~5.1 of [Sto17], let $d = \gcd(a, f)$, so that $a = d a_1$ and $f = d f_1$, let
$(u, v, w)$ be the solution of the congruences~(6) there with the degree bounds~(7) and~$w$ monic
of minimal degree, and let $d_1 = \gcd(u, w)$. In the generic case $d_1 = 1$, the
polynomials~$u$, $dv$, $w$, the polynomial~$r$ of the proposition and $\lambda = (-c)^{\deg a}$
satisfy (H1) to~(H3), with~$dv$ in the role of~$v$: (H1) is relation~(8) multiplied by~$d$, (H2)
is the second congruence in~(6) modulo~$a_1$ and holds modulo~$d$ since $d \mid b$, and~(H3)
defines~$r$. Also, $u$ is coprime to~$a$: a common prime divisor~$\pi$ of~$u$ and~$a$
divides~$dv$ by~(H2), hence~$\pi^2$ divides $\lambda a w^2$ by~(H1), so $\pi \mid w$, since~$a$ is
squarefree; this contradicts $d_1 = 1$. So cofactors as in~(H4) exist.

\begin{remark}[Formalization]
\texttt{HalvingCertificate f a b w r} is a structure with the data fields \texttt{u}, \texttt{v},
\texttt{e}, \texttt{e'} and~\texttt{lam} and one field for each condition, including
$\lambda \ne 0$; it needs only a field~$K$. It is a structure because its data come from a
computation and each condition is checked separately. The field for~(H1) is the norm identity
$v^2 - u^2 f = \lambda (a w^2)$, in the form that Lemma~\ref{lem:hv-principal} consumes; (H2)
and~(H3) are the divisibilities $a \mid v - ub$ and $w \mid ur + v$, which a concrete check proves
by supplying the quotients; (H4) is the coprimality of~$u$ and~$aw^2$ with its witnesses as data
fields. There are no monicity, degree or squarefreeness conditions: Theorem~\ref{thm:hv-generic}
needs none of them. They matter for the existence of a certificate (Proposition~5.1 of [Sto17])
and for reading $[w, r]$ as the point with the Mumford representation $(w, r)$
(Lemma~\ref{lem:mumford-exists}).
\end{remark}

\begin{lemma}[The pairs of a certificate]
\label{lem:hv-cert-mumford}
\lean{ChabautyColeman.HalvingCertificate.dvd_mul_sub_one,
  ChabautyColeman.HalvingCertificate.dvd_mul_sq_sub,
  ChabautyColeman.HalvingCertificate.dvd_mul_sub,
  ChabautyColeman.HalvingCertificate.dvd_mul_add_half,
  ChabautyColeman.HalvingCertificate.dvd_sq_sub,
  ChabautyColeman.HalvingCertificate.dvd_sq_sub_half}
\leanok
\uses{def:hv-cert}
Let $(u, v, e, e', \lambda)$ be a halving certificate for $(a, b)$ with half $(w, r)$. Then
$eu \equiv 1 \bmod a w^2$ and
\[ a w^2 \mid (ev)^2 - f\,, \qquad ev \equiv b \bmod a\,, \qquad ev \equiv -r \bmod w\,; \]
in particular $a \mid b^2 - f$ and $w \mid r^2 - f$.
\end{lemma}

\begin{proof}
\leanok
\uses{lem:hv-principal}
The first claim is~(H4). By~(H1), $v^2 - u^2 f = \lambda (a w^2)$, so
$a w^2 \mid (ev)^2 - f$ by Lemma~\ref{lem:hv-principal}\,(2). Since~$a$ and~$w$ divide~$eu - 1$,
the identities $ev - b = e(v - ub) + (eu - 1)b$ and $ev + r = e(ur + v) - (eu - 1)r$ give the two
congruences, by~(H2) and~(H3). Finally,
$b^2 - f = \bigl((ev)^2 - f\bigr) - (ev - b)(ev + b)$ and
$r^2 - f = \bigl((ev)^2 - f\bigr) - (ev + r)(ev - r)$.

\emph{Formalization.} In the namespace \texttt{HalvingCertificate}, the six statements are
separate declarations, as the proof of Theorem~\ref{thm:hv-general} uses one of them as well; the
suffix \texttt{\_half} marks the statements on the half~$(w, r)$ that have a counterpart
for~$(a, b)$.
\end{proof}

\begin{theorem}[Halving, generic case, certificate form]
\label{thm:hv-generic}
\lean{ChabautyColeman.HalvingCertificate.mumfordClass_eq_sq}
\leanok
\uses{def:hv-cert, def:hv-class}
If $(a, b)$ has a halving certificate with half $(w, r)$, then $[a, b] = 2[w, r]$ in~$J(K)$. This
is the generic case of Proposition~5.1 of [Sto17].
\end{theorem}

\begin{proof}
\leanok
\uses{lem:hv-principal, lem:hv-cert-mumford, lem:hv-factor, lem:hv-class-basic}
Let $(u, v, e, e', \lambda)$ be the certificate. By Lemma~\ref{lem:hv-principal}\,(2), applied as
in the proof of Lemma~\ref{lem:hv-cert-mumford}, $[a w^2, ev] = 0$. By
Lemma~\ref{lem:hv-cert-mumford}, $a w^2 \mid (ev)^2 - f$, $ev \equiv b \bmod a$ and
$ev \equiv -r \bmod w$; so Lemma~\ref{lem:hv-factor}, applied to~$a$ and~$w^2$ and then to~$w$
and~$w$, gives
\[ 0 = [a w^2, ev] = [a, ev] + 2[w, ev] = [a, b] + 2[w, -r]\,. \]
Since $w \mid r^2 - f$ (Lemma~\ref{lem:hv-cert-mumford}), $[w, -r] = -[w, r]$ by
Lemma~\ref{lem:hv-class-basic}\,(2). Hence $[a, b] = 2[w, r]$.

\emph{Formalization.} The conclusion reads $[a, b] = [w, r]^2$ in the multiplicative group
\texttt{HypJac K f}. Passed as terms, the field for~(H1) and \texttt{dvd\_mul\_sub\_one} determine
all implicit arguments of \texttt{mumfordClass\_eq\_one\_of\_sq\_sub\_eq}; with that lemma and those
of Lemma~\ref{lem:hv-cert-mumford}, the proof is a chain of rewrites.
\end{proof}

At the level of ideals, the same argument (Lemmas~\ref{lem:hv-principal}\,(1),
\ref{lem:hv-cert-mumford} and~\ref{lem:hv-factor}\,(1), and the ideal identities in the proof of
Lemma~\ref{lem:hv-class-basic}\,(1)) gives
$(uy - v) = \mathfrak a_{a,b}\,\mathfrak a_{w,-r}^2$; this is the identity
$\divisor(uy - v) = (a, b) + 2(w, -r) - (\deg a + 2\deg w)\infty$ in the proof of
Proposition~5.1 of [Sto17], where $(a, b)$ and~$(w, -r)$ denote divisors, as there. Neither the
degree bounds~(7) there nor the minimality of~$w$ is used.

\begin{remark}[Halves of sums]
\label{rem:hv-sums}
Definition~\ref{def:hv-cert} puts no condition on the degree of~$a$, and Proposition~5.1
of [Sto17] admits $\deg a \le 2g + 1$. So Theorem~\ref{thm:hv-generic} also certifies the halves
of a sum that is never formed. Let $P = [a_P, b_P]$ and $S = [a_S, b_S]$, put $A = a_P a_S$, and
let~$b_A$ satisfy $A \mid b_A^2 - f$, $b_A \equiv b_P \bmod a_P$ and $b_A \equiv b_S \bmod a_S$.
Then $[A, b_A] = P + S$ by Lemma~\ref{lem:hv-factor}\,(2), and a halving certificate for
$(A, b_A)$ with half $(w, r)$ shows $P + S = 2[w, r]$. This is the second variant in Remark~5.2
of [Sto17]; the first one, with $0 \ne S \in J(K)[2]$, has $\gcd(A, f) \ne 1$ and is covered by
Theorem~\ref{thm:hv-generic} if $d_1 = 1$ and by Theorem~\ref{thm:hv-general} in general.
\end{remark}

In general, $d_1 = \gcd(u, w)$ need not be~$1$ in Proposition~5.1 of [Sto17]. There
$d_1 = d_f d_a$ with $d_f = \gcd(d_1, f)$ and $d_a = \gcd(d_1, a_1)$, and
\[ u_1 = u/d_1\,, \qquad v_1 = v/d_1\,, \qquad w_1 = w/d_1\,; \]
the proposition chooses~$r$ with $r u_1 \equiv -v_1 d \bmod w_1 d_a$ and $r \equiv 0 \bmod d_f$.
Let~$R$ be the remainder of~$r$ modulo~$w_1$, and let~$\tilde b$ satisfy
$\tilde b \equiv b \bmod a/d_a$ and $\tilde b \equiv -b \bmod d_a$. Then~$u_1$, $d v_1$
and~$\lambda$, with B\'ezout cofactors for~$u_1$ and~$a w_1^2$, form a halving certificate for
$(a, \tilde b)$ with half $(w_1, R)$: (H1) is relation~(8) divided by~$d_1^2$ and multiplied
by~$d$, and~(H2) and~(H3) follow from the congruences in the proof of the proposition.
Lemmas~\ref{lem:hv-twist} and~\ref{lem:hv-split}, the content of the divisor computation at the
end of that proof, transfer $[a, \tilde b] = 2[w_1, R]$ to $[a, b] = 2[w, r]$.

\begin{theorem}[Halving in general, certificate form]
\label{thm:hv-general}
\lean{ChabautyColeman.mumfordClass_eq_sq_of_halvingCertificate}
\leanok
\uses{def:hv-cert, def:hv-class}
Let $a, b, w, r \in K[x]$ with $a \mid b^2 - f$. Assume that there are polynomials~$a'$, $d_a$,
$d_f$, $\tilde b$, $w_1$ and~$R$ in~$K[x]$ such that
\begin{enumerate}
  \item $a = a' d_a$, $\tilde b \equiv b \bmod a'$ and $\tilde b \equiv -b \bmod d_a$;
  \item $w = d_f d_a w_1$, $w \mid r^2 - f$, $r \equiv 0 \bmod d_f$, $r \equiv b \bmod d_a$
    and $r \equiv R \bmod w_1$;
  \item $(a, \tilde b)$ has a halving certificate with half $(w_1, R)$.
\end{enumerate}
Then $[a, b] = 2[w, r]$. This is Proposition~5.1 of [Sto17] in certificate form.
\end{theorem}

\begin{proof}
\leanok
\uses{thm:hv-generic, lem:hv-twist, lem:hv-factor, lem:hv-class-basic}
Lemma~\ref{lem:hv-twist} gives $[a, b] = [a, \tilde b] + 2[d_a, b]$. By
Theorem~\ref{thm:hv-generic}, $[a, \tilde b] = 2[w_1, R]$. By Lemma~\ref{lem:hv-factor}, applied
to $w = (d_f d_a) w_1$ and then to~$d_f d_a$, $2[w, r] = 2[d_f, 0] + 2[d_a, b] + 2[w_1, R]$, and
$2[d_f, 0] = 0$ by Lemma~\ref{lem:hv-class-basic}\,(2), since~$d_f$ divides~$r$ and~$r^2 - f$,
hence~$f$. So $2[w, r] = 2[d_a, b] + 2[w_1, R] = [a, b]$.

\emph{Formalization.} The certificate in~(3) is a hypothesis of type
\texttt{HalvingCertificate f a Bt W R}, with \texttt{Bt} for~$\tilde b$ and~\texttt{W}
for~$w_1$, and each condition in~(1) and~(2) is a hypothesis of its own; $a'$, $d_a$, $d_f$,
$\tilde b$, $w_1$ and~$R$ are implicit, since the hypotheses determine them. The only coprimality
used anywhere is~(H4): neither~$a'$ and~$d_a$ nor the three factors of~$w$ need to be coprime.
\end{proof}

If $d_1 = 1$, then $d_a = d_f = 1$, and the theorem with $a' = a$, $\tilde b = b$, $w_1 = w$ and
$R = r$ is Theorem~\ref{thm:hv-generic}.

\section{Cantor's algorithm and halving records}
\label{sec:hv-cantor}

When the first polynomial~$A$ of a point $[A, B]$ with $A \mid B^2 - f$ is not squarefree, write
$A = a_0^2 a$ with~$a$ squarefree; then $[A, B] = 2[a_0, B] + [a, B]$
(Lemma~\ref{lem:hv-factor}\,(3)), and a half of~$[A, B]$ is the sum of~$[a_0, B]$ and a
half~$[w, r]$ of~$[a, B]$. To obtain it as a reduced pair, the sum is computed by Cantor's
algorithm: a composition followed by reduction steps until the degree is at most~$g$. This section
turns a run of that algorithm into a certificate, as Section~\ref{sec:hv-certificates} does for
the halving. The composition needs no B\'ezout cofactors, since by Lemma~\ref{lem:hv-factor}\,(2)
three divisibilities determine the class of the composed pair. Its one special case, a common
factor of the first polynomials of the two pairs on which the pairs are opposite, cancels
(Lemma~\ref{lem:hv-cancel}). Each reduction step is an instance of Lemma~\ref{lem:hv-reduce},
which determines the new first polynomial only up to a constant factor and the new second
polynomial only by a congruence, so that a certificate needs neither the normalization of
Cantor's algorithm nor its degree conditions.

\begin{definition}[Addition certificates]
\label{def:hv-add-cert}
\lean{ChabautyColeman.AdditionCertificate}
\leanok
Let $u_1, v_1, u_2, v_2, u, v \in K[x]$. An \emph{addition certificate} for $(u_1, v_1)$
and~$(u_2, v_2)$ with sum $(u, v)$ consists of polynomials $t, p_1, p_2 \in K[x]$, an integer
$n \ge 0$, polynomials $u'_0, v'_0, \dots, u'_n, v'_n \in K[x]$ and constants
$\lambda_0, \dots, \lambda_{n-1} \in K$ such that
\begin{description}
  \item[\textmd{(A1)}] $u_1 = t p_1$, $u_2 = t p_2$ and $t \mid v_1 + v_2$;
  \item[\textmd{(A2)}] $u'_0 = p_1 p_2$, $u'_0 \mid {v'_0}^2 - f$, $v'_0 \equiv v_1 \bmod p_1$
    and $v'_0 \equiv v_2 \bmod p_2$;
  \item[\textmd{(A3)}] ${v'_i}^2 - f = \lambda_i u'_i u'_{i+1}$ and
    $v'_{i+1} \equiv -v'_i \bmod u'_{i+1}$ for $0 \le i < n$; and
  \item[\textmd{(A4)}] $(u'_n, v'_n) = (u, v)$.
\end{description}
\end{definition}

Cantor's algorithm produces such a certificate for any two pairs $(u_1, v_1)$ and~$(u_2, v_2)$
with $u_1 \mid v_1^2 - f$ and $u_2 \mid v_2^2 - f$. It takes $t = \gcd(u_1, u_2, v_1 + v_2)$;
then~$p_1$, $p_2$ and $v_1 + v_2$ generate the unit ideal of~$K[x]$. Indeed, let~$\pi$ be a common
prime divisor of the three. Then~$\pi$ divides~$t$; let~$\pi^k$, with $k \ge 1$, be the exact
power of~$\pi$ that divides~$t$. Then~$\pi^{k+1}$ divides $u_1 = t p_1$ and~$u_2$, hence
$v_1^2 - f$ and $v_2^2 - f$ and so their difference $(v_1 - v_2)(v_1 + v_2)$, whereas
$v_1 + v_2$ is divisible by~$\pi^k$, but not by~$\pi^{k+1}$, as~$t$ is not. So~$\pi$ divides
$v_1 - v_2$ as well as $v_1 + v_2$, hence~$2v_1$ and, as $2 \ne 0$ in~$K$, also~$v_1$.
Then~$\pi^2$ divides~$v_1^2$ and, as $k + 1 \ge 2$, also $v_1^2 - f$, hence~$f$, which is
impossible. So Cantor's composition (Lemma~\ref{lem:cantor-correct}), applied to $(p_1, v_1)$
and~$(p_2, v_2)$, gives~$V$ with $V \equiv v_i \bmod p_i$ for $i = 1, 2$ and
$p_1 p_2 \mid V^2 - f$, and~(A2) holds for $u'_0 = p_1 p_2$ and the remainder~$v'_0$ of~$V$
modulo~$u'_0$. The reduction steps follow the paragraph after Lemma~\ref{lem:hv-reduce}, until
$\deg u'_n \le g$.

\begin{theorem}[Cantor's algorithm, certificate form]
\label{thm:hv-add}
\lean{ChabautyColeman.AdditionCertificate.mumfordClass_mul_mumfordClass}
\leanok
\uses{def:hv-add-cert, def:hv-class}
Let $u_1, v_1, u_2, v_2, u, v \in K[x]$ with $u_1 \mid v_1^2 - f$ and $u_2 \mid v_2^2 - f$. If
$(u_1, v_1)$ and~$(u_2, v_2)$ have an addition certificate with sum $(u, v)$, then
$[u_1, v_1] + [u_2, v_2] = [u, v]$.
\end{theorem}

\begin{proof}
\leanok
\uses{lem:hv-cancel, lem:hv-factor, lem:hv-reduce}
By~(A1) and Lemma~\ref{lem:hv-cancel}, the sum $[u_1, v_1] + [u_2, v_2]$ is
$[p_1, v_1] + [p_2, v_2]$. By~(A2) and Lemma~\ref{lem:hv-factor}\,(2), this is
$[p_1 p_2, v'_0] = [u'_0, v'_0]$; by~(A3) and Lemma~\ref{lem:hv-reduce}, it is $[u'_n, v'_n]$,
which is $[u, v]$ by~(A4).

\emph{Formalization.} \texttt{AdditionCertificate}, with the six polynomials as parameters, is a
structure with the data $t$, $p_1$, $p_2$ and~$n$, the sequences $\mathbb N \to K[x]$ and
$\mathbb N \to K$ of the chain, as in Lemma~\ref{lem:hv-reduce}, and one field for each condition,
those of~(A3) for all $i < n$; it needs only a field~$K$. The fields for $t \mid v_1 + v_2$,
$v'_0 \equiv v_1 \bmod p_1$ and $v'_0 \equiv v_2 \bmod p_2$ are primed (\texttt{dvd\_add'},
\texttt{dvd\_sub\_left'}, \texttt{dvd\_sub\_right'}), as Mathlib marks such fields
(\texttt{Subgroup.mul\_mem'}), so that they do not hide Mathlib's lemmas of the same names in the
namespace. The divisibilities $u_1 \mid v_1^2 - f$ and $u_2 \mid v_2^2 - f$ are hypotheses of the
theorem, not fields: in a halving record they follow from~(R1) and~(R3). A concrete certificate
writes the sequences as lists. The certificate puts no condition on degrees or leading coefficients,
so it does not show that $(u, v)$ is reduced; that is a separate check, needed only to read~$[u, v]$
as the point with the Mumford representation $(u, v)$ (Lemma~\ref{lem:mumford-exists}). Compared
with the composition of Lemma~\ref{lem:cantor-correct}, the certificate needs no B\'ezout cofactors,
which the proof of Lemma~\ref{lem:hv-factor} obtains from Lemma~\ref{lem:hv-bezout}, and it allows a
common factor~$t$.
\end{proof}

A halving computation following Proposition~5.1 of [Sto17], for a point $[A, B]$ with
$A \mid B^2 - f$, thus produces the factorization $A = a_0^2 a$, a representative~$b$ of~$B$
modulo~$a$, the data of Theorem~\ref{thm:hv-general} for a half~$[w, r]$ of~$[a, b]$, and a run of
Cantor's algorithm for $[a_0, B] + [w, r]$. The record of such a computation is the following.

\begin{definition}[Halving records]
\label{def:hv-record}
\lean{ChabautyColeman.HalvingData}
\leanok
\uses{def:hv-cert, def:hv-add-cert}
Let $A, B, U_H, V_H \in K[x]$. A \emph{halving record} for $(A, B)$ with half $(U_H, V_H)$
consists of polynomials~$a_0$, $a$, $b$, $a'$, $d_a$, $d_f$, $\tilde b$, $w_1$, $R$, $w$ and~$r$
in~$K[x]$, a halving certificate for $(a, \tilde b)$ with half $(w_1, R)$ and an addition
certificate for $(a_0, B)$ and $(w, r)$ with sum $(U_H, V_H)$ such that
\begin{description}
  \item[\textmd{(R1)}] $A = a_0^2 a$, $A \mid B^2 - f$ and $b \equiv B \bmod a$;
  \item[\textmd{(R2)}] $a = a' d_a$, $\tilde b \equiv b \bmod a'$ and
    $\tilde b \equiv -b \bmod d_a$;
  \item[\textmd{(R3)}] $w = d_f d_a w_1$, $w \mid r^2 - f$, $r \equiv 0 \bmod d_f$,
    $r \equiv b \bmod d_a$ and $r \equiv R \bmod w_1$.
\end{description}
\end{definition}

\begin{corollary}[A halving record certifies a half]
\label{cor:hv-record}
\lean{ChabautyColeman.HalvingData.mumfordClass_eq_sq}
\leanok
\uses{def:hv-record, def:hv-class}
If $(A, B)$ has a halving record with half $(U_H, V_H)$, then $[A, B] = 2[U_H, V_H]$.
\end{corollary}

\begin{proof}
\leanok
\uses{thm:hv-general, thm:hv-add, lem:hv-factor, lem:hv-class-basic}
By~(R1), $a$ divides~$A$ and hence~$B^2 - f$; as $b^2 - f = (B^2 - f) + (b - B)(b + B)$, also
$a \mid b^2 - f$. By Lemma~\ref{lem:hv-factor}\,(3) and Lemma~\ref{lem:hv-class-basic}\,(1),
$[A, B] = 2[a_0, B] + [a, B] = 2[a_0, B] + [a, b]$, and $[a, b] = 2[w, r]$ by
Theorem~\ref{thm:hv-general}. Finally, $a_0 \mid A \mid B^2 - f$ by~(R1) and $w \mid r^2 - f$
by~(R3), so $[a_0, B] + [w, r] = [U_H, V_H]$ by Theorem~\ref{thm:hv-add}, and
$[A, B] = 2\bigl([a_0, B] + [w, r]\bigr) = 2[U_H, V_H]$.

\emph{Formalization.} \texttt{HalvingData f A B UH VH} is a structure with the data fields of the
definition (\texttt{Bt} for~$\tilde b$, \texttt{W} for~$w_1$ and \texttt{aq} for~$a'$), the
fields \texttt{cert}, a \texttt{HalvingCertificate}, and \texttt{add}, an
\texttt{AdditionCertificate}, and one field for each condition (the one for
$b \equiv B \bmod a$ is \texttt{dvd\_sub'}, primed for the reason given after
Theorem~\ref{thm:hv-add}). As in a halving certificate, the half is a parameter, so the
conclusion $[A, B] = [U_H, V_H]^2$ mentions no field of the record. If moreover $(U_H, V_H)$ is
reduced, a separate check, then $[U_H, V_H]$ is the point with the Mumford representation
$(U_H, V_H)$ (Lemma~\ref{lem:mumford-exists}).
\end{proof}

If $d_1 = 1$, then $d_a = d_f = 1$, $a' = a$, $\tilde b = b$, $w_1 = w$ and $R = r$. If~$A$ is
squarefree, then $a_0 = 1$ and $a = A$, and the addition certificate is trivial: $t = p_1 = 1$,
$p_2 = w$, $n = 0$ and $(U_H, V_H) = (w, r)$.

\section{The kernel of the descent map}
\label{sec:ds-kernel}

The inclusion $2J(K) \subseteq \ker\mu$ is formal, as~$L^\square$ has exponent~$2$. The converse
is the existence half of Proposition~5.1 of [Sto17]: if the class of~$m(a)$ is trivial, the linear
system of the proposition has a solution, and a descent ending in a halving certificate
(Theorem~\ref{thm:hv-generic} with Lemma~\ref{lem:hv-twist}) shows $[a, b] \in 2J(K)$.

\begin{theorem}[Existence of halves, {[Sto17, Proposition~5.1]}]
\label{thm:ds-halving-exists}
\lean{ChabautyColeman.exists_mumfordClass_eq_sq_of_isSquare_mumfordValue,
  ChabautyColeman.exists_sq_mul_eq_mul_sq_sub_mul_sq,
  ChabautyColeman.exists_mumfordClass_eq_sq_of_sq_sub_eq}
\leanok
\uses{def:ds-value, def:hv-class}
Let $a, b, s \in K[x]$ with~$a$ monic and squarefree, $\deg a \le 2g + 1$, $a \mid b^2 - f$ and
$s^2 \equiv (-c)^{\deg a}(a - a_1f_1) \bmod f$. Then $[a, b] \in 2J(K)$.
\end{theorem}

\begin{proof}
\leanok
\uses{thm:hv-generic, lem:hv-twist, lem:ds-value}
As in [Sto17], the polynomials~$d$, $a_1$ and~$f_1$ are squarefree and pairwise coprime, and~$b$
is prime to~$a_1$ (Lemma~\ref{lem:ds-value}\,(1)). Let $\lambda = (-c)^{\deg a}$.
\begin{enumerate}
  \item \emph{A solution.} The congruences~(6) of [Sto17],
    \[ vd \equiv ws \bmod f_1\,, \qquad vd \equiv ub \bmod a_1\,, \qquad
       uf_1 \equiv ws \bmod d\,, \]
    are linear in the coefficients of~$u$, $v$ and~$w$. Under the degree bounds~(7),
    $\deg u < \deg a/2$, $\deg v \le g + \deg a/2 - \deg d$ and $\deg w \le g$, there are
    $\deg f_1 + \deg a_1 + \deg d + 1$ coefficients and $\deg f_1 + \deg a_1 + \deg d$ linear
    conditions, so there is a nonzero solution.
  \item \emph{$w \ne 0$.} If $w = 0$, then $f_1 \mid v$, so $v = 0$ by the degree bounds (here
    $\deg a \le 2g + 1$ is used); then $a_1 \mid u$ and $d \mid u$, so $a \mid u$, and $u = 0$ by
    the degree bound. So $w \ne 0$, and~$w$ may be taken monic.
  \item \emph{The relation~(8)} $u^2f_1 = dv^2 - \lambda a_1w^2$ holds modulo~$f_1$, $a_1$ and~$d$
    by the three congruences and $s^2 \equiv \lambda(a - a_1f_1) \bmod f$, hence
    modulo~$f_1a_1d$, and all its terms have degree less than $\deg(f_1a_1d)$.
  \item \emph{The function $Uy - V$.} Let $U = u$, $V = vd$ and $W = w$. By~(8),
    $V^2 - U^2f = d(dv^2 - u^2f_1) = \lambda aW^2$, and $V \equiv Ub \bmod a$: modulo~$a_1$ by the
    second congruence, and modulo~$d$ as~$d$ divides~$b$ and~$V$.
  \item \emph{Descent.} We show by induction on $\deg W$ that $[a, b] \in 2J(K)$ whenever
    $a \mid b^2 - f$ and $U, V, W \in K[x]$ satisfy $W \ne 0$, $V^2 - U^2f = \lambda aW^2$ and
    $V \equiv Ub \bmod a$. If~$U$ is prime to~$W$, it is prime to~$aW^2$: a prime~$p$
    dividing~$U$ and~$a$ divides~$V$, so $p^2 \mid \lambda aW^2$, and $p \mid W$ as~$a$ is
    squarefree. Then~$U$, $V$, $\lambda$ and B\'ezout cofactors $e, e'$ with $eU + e'aW^2 = 1$
    form a halving certificate for $(a, b)$ with half $(W, R)$, $R = -eV$, and
    Theorem~\ref{thm:hv-generic} gives $[a, b] = 2[W, R]$. Otherwise let~$p$ be a common prime
    factor of~$U$ and~$W$. Then $p^2 \mid V^2$, so $p \mid V$, and $(U/p, V/p, W/p)$ satisfies the
    relation, with $a \mid p(V/p - (U/p)b)$. If $p \nmid a$ or $p \mid V/p - (U/p)b$, then
    $V/p \equiv (U/p)b \bmod a$, and the induction hypothesis applies. Otherwise $p \mid a$, and
    $p \mid V/p + (U/p)b$, as~$p$ divides $(V/p)^2 - (U/p)^2f$ and $b^2 - f$. Choose~$b'$ with
    $b' \equiv b \bmod a/p$ and $b' \equiv -b \bmod p$; then $a \mid b'^2 - f$,
    $V/p \equiv (U/p)b' \bmod a$, and $[a, b] = [a, b'] + 2[p, b]$ by Lemma~\ref{lem:hv-twist};
    the induction hypothesis applies to $(a, b')$.
\end{enumerate}

\emph{Formalization.} The statement takes the hypothesis that~$m(a)$ is a square in~$L$ in place
of the congruence for~$s$ (\texttt{AdjoinRoot.isSquare\_mk\_iff}), and it does not assume~$a$
monic: for general~$a$, the factor $\operatorname{lc}(a)^{-1}$ of~$m(a)$ enters~$\lambda$, which
becomes $(-c)^{\deg a}\operatorname{lc}(a)^{-1}$. Steps (1) to~(3) are
\texttt{exists\_sq\_mul\_eq\_mul\_sq\_sub\_mul\_sq}: the solution space is a product of spaces of
polynomials of bounded degree (\texttt{Polynomial.degreeLT}), the conditions are a $K$-linear map
to $K[x]/(f_1) \times K[x]/(a_1) \times K[x]/(d)$, whose dimension is
$\deg f_1 + \deg a_1 + \deg d$ (Mathlib's \texttt{AdjoinRoot.powerBasis}), and a linear map to a
space of smaller dimension has a nonzero kernel
(\texttt{Polynomial.exists\_ne\_zero\_forall\_mem\_degreeLT\_and\_dvd\_vecMul}). Step~(2) is a
parity argument there: for $w = 0$, (8) reads $u^2f_1 = dv^2$, and as $\deg(df_1)$ is odd,
$u = v = 0$ (\texttt{Polynomial.eq\_zero\_of\_mul\_sq\_eq\_mul\_sq}). Step~(5) is
\texttt{exists\_mumfordClass\_eq\_sq\_of\_sq\_sub\_eq}, by strong induction on the degree of~$W$;
it replaces the minimal solution and the factors $d_f$, $d_a$ of [Sto17], and needs only~(8) and
the second congruence. The conclusion is that $[a, b]$ is a square in the multiplicative group
\texttt{HypJac K f}.
\end{proof}

\begin{theorem}[The kernel of~$\mu$]
\label{thm:ds-kernel}
\lean{ChabautyColeman.HypJac.ker_descentMap}
\leanok
\uses{def:ds-mu}
The kernel of~$\mu$ is~$2J(K)$.
\end{theorem}

\begin{proof}
\leanok
\uses{thm:ds-halving-exists, thm:ds-formula, lem:mumford-reduce, lem:hv-factor}
$\mu(2Q) = \mu(Q)^2 = 1$. Conversely, let $\mu(P) = 1$, and let $P = [a, b]$ with a reduced pair
(Lemma~\ref{lem:mumford-reduce}): $a$ is monic of degree at most~$g$. Write $a = a_0^2a'$
with~$a'$ monic and squarefree. By Lemma~\ref{lem:hv-factor}\,(3), $P = 2[a_0, b] + [a', b]$, so
$\mu([a', b]) = 1$, and by Theorem~\ref{thm:ds-formula} there is $s \in K[x]$ with
$s(\theta)^2 = m(a') = (-c)^{\deg a'}\bigl(a'(\theta) - a'_1(\theta)f'_1(\theta)\bigr)$, that is,
$s^2 \equiv (-c)^{\deg a'}(a' - a'_1f'_1) \bmod f$, where $a' = d'a'_1$ and $f = d'f'_1$ with
$d' = \gcd(a', f)$. As $\deg a' \le g \le 2g + 1$ and $a' \mid b^2 - f$,
Theorem~\ref{thm:ds-halving-exists} gives $[a', b] \in 2J(K)$, hence $P \in 2J(K)$.

\emph{Formalization.} The kernel of \texttt{HypJac.descentMap f} is the range of
\texttt{powMonoidHom 2}. The decomposition $a = a_0^2a'$ is Mathlib's
\texttt{exists\_sq\_mul\_squarefree}; $a'$ need not be made monic, as the Lean form of
Theorem~\ref{thm:ds-halving-exists} does not assume it. The steps are
\texttt{exists\_mumfordClass\_eq\_sq\_of\_descentMap\_eq\_one} and its squarefree case
\texttt{exists\_mumfordClass\_eq\_sq\_of\_squarefree\_of\_descentMap\_eq\_one}.
\end{proof}

\begin{corollary}[The injection of $J(K)/2J(K)$]
\label{cor:ds-c0}
\lean{ChabautyColeman.HypJac.ker_descentMap_le_sup}
\leanok
\uses{def:ds-mu}
The map~$\mu$ induces an injective homomorphism $J(K)/2J(K) \to L^\square$. In particular,
$\ker\mu \subseteq \Gamma + 2J(K)$ for every subgroup $\Gamma \subseteq J(K)$.
\end{corollary}

\begin{proof}
\leanok
\uses{thm:ds-kernel}
$\ker\mu = 2J(K) \subseteq \Gamma + 2J(K)$, by Theorem~\ref{thm:ds-kernel}.

\emph{Formalization.} The Lean statement is the second claim: the kernel of
\texttt{HypJac.descentMap f} lies in the supremum of~$\Gamma$ and the range of
\texttt{powMonoidHom 2}. Restricting the codomain of~$\mu$ to a subgroup that contains its image
does not change the kernel.
\end{proof}

In the proof of Theorem~2.1 of [Sto17], the map $A(k) \to \Sel_p A$ to the $p$-Selmer group of an
abelian variety~$A$ plays the role of~$\mu$, and the inclusion of Corollary~\ref{cor:ds-c0} for it
follows from the injectivity of $A(k)/pA(k) \to \Sel_p A$. For odd-degree hyperelliptic curves,
Algorithm~4.1 and Remark~4.2 of [Sto17] use~$\mu$ in this role, with $p = 2$.

\section{Norm, \texorpdfstring{$2$}{2}-torsion, base change and rank}
\label{sec:ds-basic}

\begin{lemma}[Norms and resultants]
\label{lem:ds-resultant}
\lean{AdjoinRoot.leadingCoeff_pow_mul_norm_mk_eq_resultant}
\leanok
For every $p \in K[x]$, $c^{\deg p}N_{L/K}\bigl(p(\theta)\bigr) = \operatorname{Res}(f, p)$.
\end{lemma}

\begin{proof}
\leanok
$L = K[x]/(c^{-1}f)$, and $c^{-1}f$ is monic, so the formula for monic moduli gives
$N_{L/K}(p(\theta)) = \operatorname{Res}(c^{-1}f, p) = c^{-\deg p}\operatorname{Res}(f, p)$, by
the homogeneity of the resultant in its first argument.

\emph{Formalization.} The formula for monic moduli is \texttt{AdjoinRoot.norm\_mk\_eq\_resultant}.
As~$f$ and $c^{-1}f$ are associated, the two quotients are isomorphic by Mathlib's
\texttt{AdjoinRoot.algEquivOfAssociated}, but they are different types; the lemma transports the
norm along this isomorphism (\texttt{Algebra.norm\_eq\_of\_algEquiv}) and compares the resultants
by \texttt{Polynomial.resultant\_C\_mul\_left}. The Lean statement holds over every commutative
ring, for~$f$ with invertible leading coefficient.
\end{proof}

\begin{theorem}[The norm condition]
\label{thm:ds-norm}
\lean{ChabautyColeman.HypJac.descentMap_mem_ker_norm}
\leanok
\uses{def:ds-mu}
$N(\mu(P)) = 1$ for every $P \in J(K)$.
\end{theorem}

\begin{proof}
\leanok
\uses{thm:ds-formula, lem:ds-move, lem:ds-resultant}
Let $P = [a, b]$ with~$a$ prime to~$f$ (Lemma~\ref{lem:ds-move}). By
Lemma~\ref{lem:ds-resultant} and the symmetry of the resultant,
\[ N_{L/K}\bigl(m(a)\bigr) = \bigl((-c)^{\deg a}\operatorname{lc}(a)^{-1}\bigr)^{2g+1}
   c^{-\deg a}(-1)^{(2g+1)\deg a}\operatorname{Res}(a, f)\,. \]
Over a splitting field of~$a$, $\operatorname{Res}(a, f) =
\operatorname{lc}(a)^{2g+1}\prod_\alpha f(\alpha) = \operatorname{lc}(a)^{2g+1}\bigl(\prod_\alpha
b(\alpha)\bigr)^2$, the products running over the roots~$\alpha$ of~$a$ with multiplicity, since
$f(\alpha) = b(\alpha)^2$; and $\prod_\alpha b(\alpha) = \operatorname{lc}(a)^{-\deg
b}\operatorname{Res}(a, b) \in K$. Collecting the factors, $N_{L/K}(m(a)) = \bigl(c^{g\deg
a}\operatorname{lc}(a)^{-\deg b}\operatorname{Res}(a, b)\bigr)^2$.

\emph{Formalization.} The statement is that~$\mu(P)$ lies in the kernel of the map
\texttt{Units.ModPow.map (Algebra.norm K) 2} induced by the norm. The proof needs no splitting
field: for $a \mid b^2 - f$, the resultant $\operatorname{Res}(a, f)$ is computed in an even
degree from $f \equiv b^2 \bmod a$
(\texttt{Polynomial.isSquare\_leadingCoeff\_pow\_mul\_resultant\_of\_dvd\_sq\_sub}), and
\texttt{isSquare\_norm\_mumfordValue} collects the constants.
\end{proof}

\begin{theorem}[$2$-torsion]
\label{thm:ds-two-torsion}
\lean{ChabautyColeman.HypJac.sq_eq_one_iff, ChabautyColeman.mumfordClass_zero_eq_of_mul_eq,
  ChabautyColeman.mumfordClass_zero_inj, ChabautyColeman.HypJac.card_ker_powMonoidHom_two}
\leanok
\uses{def:hv-class}
\begin{enumerate}
  \item A point $P \in J(K)$ satisfies $2P = 0$ if and only if $P = [h, 0]$ for a divisor~$h$
    of~$f$.
  \item $[h, 0] = [f/h, 0]$ for every nonzero divisor~$h$ of~$f$, and the points $[h, 0]$ with
    $h \mid f$ monic and $\deg h \le g$ are pairwise distinct.
  \item $J(K)[2]$ has~$2^{m-1}$ elements, where~$m$ is the number of irreducible factors of~$f$.
\end{enumerate}
\end{theorem}

\begin{proof}
\leanok
\uses{lem:hv-class-basic, lem:mumford-reduce, lem:mumford-exists, lem:hv-reduce}
(1) If $h \mid f$, then $2[h, 0] = 0$ by Lemma~\ref{lem:hv-class-basic}\,(2). Conversely, let
$P = [u, v]$ with a reduced pair (Lemma~\ref{lem:mumford-reduce}). Then $(u, -v)$ is reduced as
well, and $[u, -v] = -P$; so $2P = 0$ gives $v = -v$ by the uniqueness of reduced representatives
(Lemma~\ref{lem:mumford-exists}), hence $v = 0$ and $u \mid -f$. (2) As
$0^2 - f = (-1) \cdot h \cdot (f/h)$, Lemma~\ref{lem:hv-reduce} with $\lambda = -1$ gives
$[h, 0] = [f/h, 0]$. A pair $(h, 0)$ with~$h$ monic and $\deg h \le g$ is reduced, so these points
are distinct by Lemma~\ref{lem:mumford-exists}. (3) The~$2^m$ monic divisors~$h$ of~$f$ fall into
pairs $\{h, c^{-1}f/h\}$, of which exactly one member has degree at most~$g$, as the degrees add
up to $2g + 1$; by~(1) and~(2), the points of order dividing~$2$ are the $2^{m-1}$ distinct points
$[h, 0]$ with $\deg h \le g$.

\emph{Formalization.} Mathlib has no subgroup of $n$-torsion points of a multiplicative group, so
$J(K)[2]$ is the kernel of the squaring map \texttt{powMonoidHom 2}, and~$m$ is the cardinality of
the finite set \texttt{f.Factors} of monic irreducible factors. In~(3), a monic divisor of~$f$ is a
product over a subset of \texttt{f.Factors} (\texttt{Polynomial.Factors.exists\_prod\_eq}), a subset
and its complement give the pair $\{h, c^{-1}f/h\}$, and
\texttt{Finset.natCard\_subtype\_eq\_two\_pow\_of\_compl\_iff\_not} counts the subsets with
$2\deg h < \deg f$. That $[h, 0]$ has order dividing~$2$ is \texttt{mumfordClass\_zero\_sq}.
In~(2), \texttt{mumfordClass\_zero\_eq\_of\_mul\_eq} takes the cofactor~$e$ with $he = f$ and is
one Cantor step, and \texttt{mumfordClass\_zero\_inj} is the distinctness for monic~$h$ with
$2\deg h < \deg f$, from \texttt{mumfordClass\_inj}, the form for classes of
\texttt{mumfordIdeal\_reduced\_unique}, which also gives~(1).
\end{proof}

\begin{theorem}[Base change]
\label{thm:ds-naturality}
\lean{ChabautyColeman.HypJac.map_descentMap, ChabautyColeman.HypJac.map_comp_descentMap}
\leanok
\uses{def:ds-mu, def:basechange}
Let $K \to K'$ be a field extension over which~$f$ stays squarefree (for instance, $K$ perfect),
let $L' = K'[x]/(f) = L \otimes_K K'$, and let $\mu' \colon J(K') \to L'^\square$ be the map of
Definition~\ref{def:ds-mu} over~$K'$. Then $\mu'(P_{K'}) = \mu(P)_{L'}$ for every $P \in J(K)$,
where $P \mapsto P_{K'}$ is the map $J(K) \to J(K')$ of Definition~\ref{def:basechange} and
$\xi \mapsto \xi_{L'}$ the map $L^\square \to L'^\square$.
\end{theorem}

\begin{proof}
\leanok
\uses{thm:ds-formula}
Let $P = [a, b]$ with~$a$ prime to~$f$. Then $P_{K'} = [a, b]$ over~$K'$, $a$ is prime to~$f$
over~$K'$ by the same B\'ezout identity, and~$m(a)$ computed in~$L'$ is the image of~$m(a)$
computed in~$L$.

\emph{Formalization.} $L'$ is \texttt{AdjoinRoot (f.map (algebraMap K K'))}, the map $L \to L'$
is \texttt{AdjoinRoot.map}, and on square classes \texttt{Units.ModPow.map}; the map on Jacobians
is \texttt{HypJac.map}, and \texttt{HypJac.map\_mumfordClass} computes it on Mumford classes.
\texttt{HypJac.map\_descentMap} is the pointwise statement and \texttt{HypJac.map\_comp\_descentMap}
the square of homomorphisms; as~$m(a)$ commutes with every homomorphism of fields
(\texttt{map\_mumfordValue}), the proof uses Theorem~\ref{thm:ds-formula} on any representative,
not one prime to~$f$. The hypotheses are the instances
\texttt{Fact (Squarefree (f.map (algebraMap K K')))}, which
\texttt{instFactSquarefreeMapAlgebraMap} provides for a perfect field~$K$, and
\texttt{NeZero (2 : K')}, which holds for every extension of~$K$.
\end{proof}

For a number field~$K$ and a place~$v$, with $K' = K_v$, the theorem says that the local descent
map~$\mu_v$ takes the image of~$P$ in~$J(K_v)$ to the restriction of~$\mu(P)$ to $L \otimes_K
K_v$. The map $J(K) \to J(K')$ is injective (Theorem~\ref{thm:basechange-inj}).

\begin{theorem}[Rank]
\label{thm:ds-rank}
\lean{ChabautyColeman.HypJac.card_range_descentMap,
  ChabautyColeman.HypJac.pow_rank_le_card_of_range_descentMap_le}
\leanok
\uses{def:ds-mu}
Assume that~$J(K)$ is finitely generated, of rank~$\operatorname{rk} J(K)$. Then
\[ \#\mu(J(K)) = 2^{\operatorname{rk} J(K)}\,\#J(K)[2]\,. \]
Consequently $2^{\operatorname{rk} J(K)}\,\#J(K)[2] \le \#S$ for every finite subgroup
$S \subseteq L^\square$ containing~$\mu(J(K))$.
\end{theorem}

\begin{proof}
\leanok
\uses{thm:ds-kernel}
$\mu(J(K)) \cong J(K)/2J(K)$ by Theorem~\ref{thm:ds-kernel}, and a finitely generated abelian
group $G \cong \ZZ^n \times T$ with~$T$ finite has $\#G/2G = 2^n\,\#T/2T = 2^n\,\#G[2]$.

\emph{Formalization.} The hypothesis is the instance \texttt{IsMulFG (HypJac K f)}, the rank is
\texttt{Module.finrank} over~$\ZZ$ of the additive group \texttt{Additive (HypJac K f)}, and the
count is \texttt{CommGroup.card\_modPow}. For finitely generated~$J(K)$ this rank is
$\dim_{\QQ}(J(K) \otimes_{\ZZ} \QQ)$, the rank of Definition~\ref{def:rank-hyp} in the hypothesis of
the main theorems; that definition assumes no finite generation, and the comparison of the two
forms is not formalized.
\end{proof}

\section{Over a Dedekind domain}
\label{sec:ds-dedekind}

\begin{setting}
\label{set:ds-dedekind}
In Setting~\ref{set:ds-setting}, let~$K$ be a perfect field (Remark~\ref{rem:ds-perfect}), the
fraction field of a Dedekind domain~$R$, and let~$\Sigma$ be a set of nonzero primes of~$R$. For
a finite extension field~$F$ of~$K$ with integral closure~$B$ of~$R$ and $n \ge 1$, let
$F(\Sigma, n) \subseteq F^\times/(F^\times)^n$ be the group of classes whose valuation is
divisible by~$n$ at every prime of~$B$ not above~$\Sigma$ (the Selmer group of~$B$ relative to the
primes above~$\Sigma$). For each monic irreducible factor~$p$ of~$f$, let~$B_p$ be the integral
closure of~$R$ in $L_p = K[x]/(p)$.
\end{setting}

\begin{definition}[The group $A(\Sigma, 2)$]
\label{def:ds-selmer}
\lean{AdjoinRoot.selmerGroup, AdjoinRoot.finite_selmerGroup}
\leanok
Let $A(\Sigma, 2) \subseteq L^\square$ be the group of classes whose component in~$L_p^\square$
lies in~$L_p(\Sigma, 2)$ for every~$p$, under $L^\square \cong \prod_p L_p^\square$. If~$\Sigma$
is finite and each~$B_p$ has a finite class group and a finitely generated unit group (for
instance if~$K$ is a number field), then $A(\Sigma, 2)$ is finite.
\end{definition}

\begin{proof}
\leanok
Each $L_p(\Sigma, 2)$ is finite, by the exact sequence
$1 \to B_{p,\Sigma}^\times/2 \to L_p(\Sigma, 2) \to \Cl(B_{p,\Sigma})[2] \to 1$ for the
ring~$B_{p,\Sigma}$ of $\Sigma$-integers and the $S$-unit theorem.

\emph{Formalization.} \texttt{AdjoinRoot.selmerGroup} is defined for every exponent~$n$ and used
with $n = 2$. It is the pullback along \texttt{AdjoinRoot.modPowEquivPiFactors} of the product,
over \texttt{f.Factors}, of the Selmer groups \texttt{IsDedekindDomain.selmerGroupAbove} of the
integral closures of~$R$ in the factors \texttt{AdjoinRoot p}, relative to the primes above~$\Sigma$
(\texttt{IsDedekindDomain.HeightOneSpectrum.primesAbove}); \texttt{AdjoinRoot.mem\_selmerGroup\_iff}
and \texttt{AdjoinRoot.mk\_mem\_selmerGroup\_iff} are the componentwise conditions. The
separability of the factor fields is an instance family, which
\texttt{AdjoinRoot.instIsSeparableFactors} provides for a perfect~$K$.
\texttt{AdjoinRoot.finite\_selmerGroup} takes a finite~$\Sigma$ and instances for the finiteness of
the class groups and the finite generation of the unit groups of the integral closures. Mathlib's
\texttt{IsDedekindDomain.selmerGroup} is the case of the fraction field itself.
\end{proof}

\begin{remark}[Separability]
\label{rem:ds-perfect}
The integral closure of~$R$ in~$L_p$ is a Dedekind domain finite over~$R$ by Mathlib's results for
finite \emph{separable} extensions, and Lemma~\ref{lem:ds-norm-selmer} assumes the extension
separable, as Mathlib's relative norm \texttt{Algebra.intNorm} does; in the proof of
Lemma~\ref{lem:ds-point-mumford} it is applied to $F[x]/(a)$ for a finite extension~$F$ of~$K$ and
an irreducible~$a \in F[x]$. Both separability conditions hold if~$K$ is perfect (then so is~$F$),
in particular in characteristic~$0$, which covers number fields and their completions. Over an
imperfect~$K$, $f$ is assumed squarefree, not separable, and an irreducible~$a$ can be
inseparable. So~$K$ is assumed perfect in this section and the next.
\end{remark}

\begin{lemma}[Norms map Selmer groups to Selmer groups]
\label{lem:ds-norm-selmer}
\lean{IsDedekindDomain.norm_mem_selmerGroup}
\leanok
Let $R \subseteq B$ be Dedekind domains with fraction fields $K \subseteq F$, where~$F/K$ is finite
and separable and~$B$ is the integral closure of~$R$ in~$F$, let~$\Sigma$ be a set of nonzero
primes of~$R$, and $n \ge 1$. If $x \in F^\times$ has valuation divisible by~$n$ at every prime
of~$B$ not above~$\Sigma$, then~$N_{F/K}(x)$ has valuation divisible by~$n$ at every prime of~$R$
not in~$\Sigma$. So~$N_{F/K}$ maps $F(\Sigma, n)$ into $K(\Sigma, n)$.
\end{lemma}

\begin{proof}
\leanok
For a prime $v \notin \Sigma$, $v(N_{F/K}x) = \sum_{w \mid v} f(w \mid v)\,w(x)$, the sum over the
primes~$w$ of~$B$ above~$v$ with their residue degrees: for $x \in B$, the ideal norm maps $xB$ to
$N_{F/K}(x)R$ and each~$w$ to $v^{f(w \mid v)}$, and is multiplicative; in general write
$x = y/z$ with $y, z \in B$.

\emph{Formalization.} The map induced by the norm sends \texttt{selmerGroupAbove} into Mathlib's
\texttt{IsDedekindDomain.selmerGroup}. The valuation formula is
\texttt{IsDedekindDomain.toAdd\_valuationOfNeZero\_norm}, from \texttt{Ideal.multiplicity\_relNorm}
(the multiplicity of~$v$ in the relative norm of an ideal is $\sum_{w \mid v} f(w \mid v)$ times
that of~$w$), which rests on Mathlib's \texttt{Ideal.relNorm}, \texttt{Ideal.relNorm\_singleton}
and \texttt{Ideal.relNorm\_eq\_pow\_of\_isMaximal}. The lemma is missing from Mathlib.
\end{proof}

\begin{lemma}[The point lemma]
\label{lem:ds-point}
\lean{AdjoinRoot.of_sub_root_mem_selmerGroup,
  IsDedekindDomain.mem_selmerGroupAbove_of_sq_eq_aeval}
\leanok
\uses{def:ds-selmer}
Assume that~$\Sigma$ contains every prime~$v$ of~$R$ at which a coefficient of~$f$ is not
integral, $c$ is not a unit or $\operatorname{disc} f$ is not a unit. Let $x_0, y_0 \in K$ with
$y_0^2 = f(x_0) \ne 0$. Then $x_0 - \theta$ is a unit of~$L$, and its class lies in
$A(\Sigma, 2)$.
\end{lemma}

\begin{proof}
\leanok
For each~$p$, $p(x_0) \ne 0$, so the component $x_0 - \theta_p$ is not zero. Let~$w$ be a prime
of~$B_p$ above a prime $v \notin \Sigma$. As $c^{-1}f$ is monic with $v$-integral coefficients,
$w(\theta_p) \ge 0$. If $w(x_0) < 0$, then $w(x_0 - \theta_p) = w(x_0)$ and
$2w(y_0) = w(f(x_0)) = (2g + 1)w(x_0)$, as the leading term of~$f(x_0)$ dominates ($c$ is a
$v$-unit); so~$w(x_0)$ is even. If $w(x_0) \ge 0$, write
$f(x_0) = f(x_0) - f(\theta_p) = (x_0 - \theta_p)G$ with~$G$ integral at~$w$ and
$G \equiv f'(\theta_p) \bmod (x_0 - \theta_p)$. As $\operatorname{disc} f$ and~$c$ are $v$-units,
so is $N_{L/K}(f'(\theta))$, and~$f'(\theta_p)$ is a $w$-unit. So either $w(x_0 - \theta_p) = 0$,
or $w(G) = 0$ and $w(x_0 - \theta_p) = w(f(x_0)) = 2w(y_0)$; in both cases $w(x_0 - \theta_p)$ is
even.

\emph{Formalization.} \texttt{AdjoinRoot.of\_sub\_root\_mem\_selmerGroup} is stated for the class
of any unit with value $x_0 - \theta$ (\texttt{AdjoinRoot.isUnit\_of\_sub\_root} supplies one). The
condition on~$\Sigma$ is that~$f$ has separable reduction at the primes not in~$\Sigma$
(\texttt{Polynomial.HasSeparableReduction}: integral coefficients, a unit leading coefficient and a
unit discriminant, Mathlib's \texttt{Polynomial.discr}). The valuation core is
\texttt{Valuation.exists\_map\_sub\_eq\_sq}, for a polynomial of odd degree with a unit leading
coefficient: for non-integral~$x_0$ the leading term dominates; otherwise a Taylor expansion over
the integers (\texttt{Polynomial.exists\_mem\_eval\_eq\_add\_mul}) and the unit~$f'(\theta_p)$
(\texttt{Valuation.map\_derivative\_eval\_eq\_one}, from a B\'ezout identity for the resultant
of~$f$ and~$f'$ over the valuation ring, in place of the norm argument) give
\texttt{Valuation.map\_sub\_eq\_one\_or\_eq\_map\_eval}. The local form
\texttt{Polynomial.HasSeparableReduction.exists\_map\_sub\_eq\_sq} is applied over the extension
field, to the image of~$f$, which keeps separable reduction
(\texttt{Polynomial.HasSeparableReduction.map}); so~$x_0$ and~$y_0$ may lie in the extension field,
as the proof of Lemma~\ref{lem:ds-point-mumford} needs, and
\texttt{IsDedekindDomain.mem\_selmerGroupAbove\_of\_sq\_eq\_aeval} is the point lemma in a single
field.
\end{proof}

\begin{lemma}[The point lemma for Mumford polynomials]
\label{lem:ds-point-mumford}
\lean{IsDedekindDomain.mem_selmerGroupAbove_of_dvd_sq_sub,
  IsDedekindDomain.mem_selmerGroup_of_dvd_sq_sub}
\leanok
Let~$\Sigma$ be as in Lemma~\ref{lem:ds-point}, let~$F$ be a finite extension field of~$K$,
$\theta \in F$ a root of~$f$, and $a, b \in F[x]$ with $a \mid b^2 - f$ and~$a$ prime to~$f$. Then
$\operatorname{lc}(a)^{-1}a(\theta) \ne 0$, and its class lies in~$F(\Sigma, 2)$.
\end{lemma}

\begin{proof}
\leanok
\uses{lem:ds-point, lem:ds-norm-selmer, lem:ds-resultant}
As~$a$ is prime to~$f$ and $f(\theta) = 0$, $a(\theta) \ne 0$. If $a = a'a''$, then~$a'$
and~$a''$ divide $b^2 - f$ and are prime to~$f$, and $\operatorname{lc}(a)^{-1}a(\theta)$ is the
product of the values for~$a'$ and~$a''$; a constant~$a$ gives~$1$. So, by induction on the
irreducible factors, we may assume that~$a$ is monic and irreducible. Let $F' = F[x]/(a)$, a
finite extension field of~$F$, and let~$x_1 \in F'$ be the image of~$x$. Then
$b(x_1)^2 = f(x_1)$, and $f(x_1) \ne 0$ as~$a$ is prime to~$f$. Let~$w$ be a prime of the integral
closure of~$R$ in~$F'$ above a prime $v \notin \Sigma$. The proof of Lemma~\ref{lem:ds-point},
with~$F'$, $\theta$, $x_1$ and~$b(x_1)$ in place of~$L_p$, $\theta_p$, $x_0$ and~$y_0$, shows that
$w(x_1 - \theta)$ is even: it uses only that~$\theta$ is a root of~$f$ and that
$y_0^2 = f(x_0) \ne 0$. So the class of $\theta - x_1$ lies in~$F'(\Sigma, 2)$, and by
Lemma~\ref{lem:ds-norm-selmer}, for $F'/F$ and the primes of the integral closure of~$R$ in~$F$
above~$\Sigma$, its norm lies in~$F(\Sigma, 2)$. As the characteristic polynomial of~$x_1$ over~$F$
is~$a$, this norm is $N_{F'/F}(\theta - x_1) = a(\theta)$; equivalently, by
Lemma~\ref{lem:ds-resultant} for~$a$ over~$F$, it is $\operatorname{Res}(a, \theta - x) =
a(\theta)$.

\emph{Formalization.} The general form is
\texttt{IsDedekindDomain.mem\_selmerGroup\_of\_dvd\_sq\_sub}, over a Dedekind domain~$B$ with
perfect fraction field (in the application the integral closure of~$R$ in~$L_p$), for a polynomial
with separable reduction outside a set of primes of~$B$: the class of a unit~$u$ with
$u = \operatorname{lc}(a)^{-1}a(\theta)$ lies in Mathlib's \texttt{IsDedekindDomain.selmerGroup}.
The induction is \texttt{WfDvdMonoid.induction\_on\_irreducible}. The step for an irreducible~$a$,
monic or not, works in $F' = \texttt{AdjoinRoot a}$ with the integral closure of~$B$ in~$F'$: the
point lemma in a field and Lemma~\ref{lem:ds-norm-selmer} give
\texttt{IsDedekindDomain.norm\_sub\_mem\_selmerGroup\_of\_sq\_eq\_aeval}, and
\texttt{AdjoinRoot.leadingCoeff\_mul\_norm\_of\_sub\_root} computes
$N_{F'/F}(\theta - x_1) = \operatorname{lc}(a)^{-1}a(\theta)$. The lemma as stated, with~$f$
over~$K$ and~$\Sigma$, is \texttt{IsDedekindDomain.mem\_selmerGroupAbove\_of\_dvd\_sq\_sub}:
separable reduction moves up to~$B$ by \texttt{Polynomial.HasSeparableReduction.map}.
\end{proof}

\begin{theorem}[The image of~$\mu$]
\label{thm:ds-image-selmer}
\lean{ChabautyColeman.HypJac.range_descentMap_le_selmerGroup}
\leanok
\uses{def:ds-mu, def:ds-selmer}
If~$\Sigma$ is as in Lemma~\ref{lem:ds-point}, then $\mu(J(K)) \subseteq A(\Sigma, 2)$.
\end{theorem}

\begin{proof}
\leanok
\uses{lem:ds-move, lem:ds-point-mumford}
Let $P \in J(K)$. By Lemma~\ref{lem:ds-move}, $P = [a, b]$ with~$a$ prime to~$f$, and by
Definition~\ref{def:ds-mu}, $\mu(P)$ is the class of
$m(a) = (-c)^{\deg a}\operatorname{lc}(a)^{-1}a(\theta)$. Its component at~$p$ is
$(-c)^{\deg a}\operatorname{lc}(a)^{-1}a(\theta_p)$. By Lemma~\ref{lem:ds-point-mumford} for
$F = L_p$ and $\theta = \theta_p$ ($a$ stays prime to~$f$ in~$L_p[x]$), the class of
$\operatorname{lc}(a)^{-1}a(\theta_p)$ lies in~$L_p(\Sigma, 2)$, and so does the class
of~$(-c)^{\deg a}$, as~$c$ is a unit at every prime not in~$\Sigma$. So every component
of~$\mu(P)$ lies in~$L_p(\Sigma, 2)$, that is, $\mu(P) \in A(\Sigma, 2)$.

\emph{Formalization.} The proof works one component~$L_p$ at a time; only norms of field extensions
$L_p[x]/(a')$ of~$L_p$ occur, in Lemma~\ref{lem:ds-point-mumford}. It combines
\texttt{HypJac.exists\_isCoprime\_mumfordClass\_eq},
\texttt{HypJac.descentMap\_mumfordClass\_of\_isCoprime} and
\texttt{AdjoinRoot.mem\_selmerGroup\_iff} with \texttt{algHom\_mumfordValue\_mem\_selmerGroupAbove}:
for a $K$-algebra homomorphism $\varphi \colon L \to F$ (here the projection to~$L_p$) the class of
$\varphi(m(a))$ lies in~$F(\Sigma, 2)$, by Lemma~\ref{lem:ds-point-mumford} for the
root~$\varphi(\theta)$ of~$f$ and
\texttt{IsDedekindDomain.mem\_selmerGroupAbove\_of\_valuation\_eq\_one} for the constant. Each~$L_p$
is perfect by \texttt{Algebra.IsAlgebraic.perfectField}.
\end{proof}

\section{Number fields}
\label{sec:ds-number}

\begin{setting}
\label{set:ds-number}
In Setting~\ref{set:ds-setting}, let~$K$ be a number field and $R = O_K$. For a place~$v$ of~$K$,
let~$K_v$ be the completion, $L_v = K_v[x]/(f) = L \otimes_K K_v$,
$\mu_v \colon J(K_v) \to L_v^\square$ the map of Definition~\ref{def:ds-mu} over~$K_v$ ($f$ stays
squarefree, as~$K$ has characteristic~$0$), $r_v \colon J(K) \to J(K_v)$ the natural map and
$\operatorname{res}_v \colon L^\square \to L_v^\square$ the restriction. Let~$\Sigma_f$ be the
finite set of primes~$v$ of~$R$ at which a coefficient of~$f$ is not integral, $c$ is not a unit
or $\operatorname{disc} f$ is not a unit.
\end{setting}

\begin{definition}[The $2$-Selmer group]
\label{def:ds-sel2}
\lean{ChabautyColeman.HypJac.selmerGroupTwo, ChabautyColeman.HypJac.localCondition}
\leanok
\uses{def:ds-mu}
The $2$-Selmer group is
\[ \Sel_2 = \{\xi \in L^\square : N(\xi) = 1 \text{ and } \operatorname{res}_v(\xi) \in
   \mu_v(J(K_v)) \text{ for every place~$v$ of~$K$}\}\,. \]
\end{definition}

\begin{theorem}[Finiteness of the $2$-Selmer group]
\label{thm:ds-sel2}
\lean{ChabautyColeman.HypJac.range_descentMap_le_selmerGroupTwo,
  ChabautyColeman.HypJac.selmerGroupTwo_le_selmerGroup,
  ChabautyColeman.HypJac.selmerGroupTwo_le_ker_norm,
  ChabautyColeman.HypJac.finite_selmerGroupTwo}
\leanok
\uses{def:ds-sel2, def:ds-selmer}
\begin{enumerate}
  \item $\mu(J(K)) \subseteq \Sel_2$.
  \item $\Sel_2 \subseteq A(\Sigma_f, 2) \cap \ker N$.
  \item $\Sel_2$ is finite.
\end{enumerate}
\end{theorem}

\begin{proof}
\leanok
\uses{thm:ds-norm, thm:ds-naturality, thm:ds-image-selmer}
(1) Theorem~\ref{thm:ds-norm}, and $\operatorname{res}_v(\mu(P)) = \mu_v(r_v(P))$ by
Theorem~\ref{thm:ds-naturality}. (2) For a prime $v \notin \Sigma_f$,
Theorem~\ref{thm:ds-image-selmer} over~$K_v$, with $R = O_v$ and $\Sigma = \emptyset$, shows
that $\mu_v(J(K_v))$ consists of unramified classes; and a class in~$L^\square$ lies in
$A(\Sigma_f, 2)$ if and only if its restriction is unramified at every prime $v \notin \Sigma_f$.
(3) follows from~(2) and Definition~\ref{def:ds-selmer}.

\emph{Formalization.} \texttt{HypJac.selmerGroupTwo} is defined for a Dedekind domain~$R$ whose
fraction field~$K$ has characteristic~$0$ and a family of field extensions of~$K$ (for a number
field, the completions at the infinite places): it is the intersection of $\ker N$, of the local
conditions at the completions \texttt{v.adicCompletion K} at the primes~$v$ of~$R$, and of the
local conditions at the members of the family. The local condition \texttt{HypJac.localCondition}
at a field extension~$F$ of~$K$ is the preimage of the image of~$\mu_F$ under the restriction
to~$F$. (2) is stated for every set~$\Sigma$ of primes outside which~$f$ has separable reduction
(\texttt{Polynomial.HasSeparableReduction}), and of the equivalence in its proof only the
implication from right to left is used:
\texttt{AdjoinRoot.mem\_selmerGroup\_of\_forall\_adicCompletion}, for every squarefree~$f$ and every
exponent. There~$O_v$ is the valuation ring of the completion, a discrete valuation ring, and~$f$
has separable reduction at its prime because it has at~$v$. In~(3), $\Sigma_f$ is finite by
\texttt{Polynomial.finite\_setOf\_not\_hasSeparableReduction}, as $\operatorname{disc} f \ne 0$, and
\texttt{HypJac.finite\_selmerGroupTwo} takes the instances on class groups and unit groups of
\texttt{AdjoinRoot.finite\_selmerGroup}.
\end{proof}

\begin{definition}[Cited structures]
\label{def:ds-cited}
\lean{ChabautyColeman.HypJac.MordellWeil, ChabautyColeman.HypJac.LocalStructure}
\leanok
\begin{description}
  \item[\textmd{(MW)}] For a number field~$K$, the group~$J(K)$ is finitely generated (the
    Mordell--Weil theorem: A.~Weil, \emph{L'arithm\'etique sur les courbes alg\'ebriques}, Acta
    Math.~52 (1929), 281--315; M.~Hindry and J.~H.~Silverman, \emph{Diophantine geometry. An
    introduction}, Springer, 2000, Theorem~C.0.1).
  \item[\textmd{(Loc)}] For a finite extension~$K_v$ of~$\QQ_\ell$ with ring of
    integers~$O_v$, there is an isomorphism of groups $J(K_v) \cong O_v^g \times T$ with~$T$
    finite (A.~Mattuck, \emph{Abelian varieties over $p$-adic ground fields}, Ann.\ of Math.~(2)
    62 (1955), 92--119; for $K_v = \QQ_2$ also B.~Poonen and M.~Stoll, \emph{Most odd degree
    hyperelliptic curves have only one rational point}, Ann.\ of Math.~(2) 180 (2014),
    1137--1166, Section~4).
\end{description}
\end{definition}

Mattuck's isomorphism is one of topological groups. The statements below use only the group
isomorphism, and no topology on~$J(K_v)$ is involved.

\begin{remark}[Formalization]
The two results are cited, not proved: each is a class or structure whose docstring gives the
references. \texttt{HypJac.MordellWeil K f}, for a number field~$K$, is a class with the single
field that \texttt{HypJac K f} is finitely generated (\texttt{IsMulFG}), which it provides as an
instance. \texttt{HypJac.LocalStructure p F f g}, for a field~$F$ that is an algebra over the
field~$\QQ_p$ of Mathlib, is a structure with values in \texttt{Prop}: $\deg f = 2g + 1$, and there
is a finite group~$T$ with a group isomorphism from \texttt{HypJac F f} to the product of
$(\ZZ_p^d)^g$ and~$T$, where $d = [F : \QQ_p]$. As a group, $O_v \cong \ZZ_p^d$; this form
of~(Loc) avoids the integral closure of~$\ZZ_p$ in~$F$.
\end{remark}

\begin{lemma}[Consequences of (Loc)]
\label{lem:ds-loc}
\lean{ChabautyColeman.HypJac.LocalStructure.card_modPow_two,
  ChabautyColeman.HypJac.LocalStructure.card_modPow_two_of_ne_two,
  ChabautyColeman.HypJac.LocalStructure.mem_iInf_range_powMonoidHom_pow_iff}
\leanok
\uses{def:ds-cited, def:hypjac}
Assume~(Loc) for a finite extension~$K_v$ of~$\QQ_\ell$.
\begin{enumerate}
  \item $\#J(K_v)/2J(K_v) = 2^{g[K_v : \QQ_2]}\,\#J(K_v)[2]$ if $\ell = 2$, and
    $\#J(K_v)/2J(K_v) = \#J(K_v)[2]$ if~$\ell$ is odd.
  \item $\bigcap_{n \ge 0} \ell^nJ(K_v)$ is the subgroup of elements of finite order prime
    to~$\ell$.
\end{enumerate}
\end{lemma}

\begin{proof}
\leanok
Let $J(K_v) \cong O_v^g \times T$. As a group, $O_v \cong \ZZ_\ell^{[K_v : \QQ_\ell]}$ is
torsion-free, so $J(K_v)[2] \cong T[2]$, and $\#T/2T = \#T[2]$ as~$T$ is finite.
(1) $\#O_v/2O_v$ is $2^{[K_v : \QQ_2]}$ if $\ell = 2$ and~$1$ if~$\ell$ is odd, as~$2$ is then a
unit. (2) In every abelian group, an element of finite order~$m$ prime to~$\ell$ is an
$\ell^n$-th multiple for every~$n$: if $k\ell^n \equiv 1 \bmod m$, then $x = \ell^n(kx)$.
Conversely, $\bigcap_n \ell^nO_v = 0$, and the subgroups $\ell^nT$ of the finite group~$T$
decrease and become constant, equal to the subgroup~$T_{\ell'}$ of elements of order prime
to~$\ell$ (on the $\ell$-primary part of~$T$, multiplication by~$\ell^n$ is zero for large~$n$, and
on~$T_{\ell'}$ it is bijective); as $\ell^n(O_v^g \times T) = \ell^nO_v^g \times \ell^nT$, the
intersection is $0 \times T_{\ell'}$.

\emph{Formalization.} With the prime~$\ell$ named~$p$, (1) is
\texttt{HypJac.LocalStructure.card\_modPow\_two} for $p = 2$ and
\texttt{HypJac.LocalStructure.card\_modPow\_two\_of\_ne\_two} for $p \ne 2$, about the cardinality
of the quotient by the range of \texttt{powMonoidHom 2} and of its kernel; both go through the
index of $2\ZZ_p$ in~$\ZZ_p$, which is~$2$ for $p = 2$ (\texttt{PadicInt.index\_range\_nsmul\_self})
and~$1$ otherwise (\texttt{PadicInt.range\_nsmul\_eq\_top\_of\_not\_dvd}), and
\texttt{AddSubgroup.index\_range\_nsmul\_prod\_of\_injective}. (2) says that an element lies in the
intersection of the ranges of the maps \texttt{powMonoidHom (p\textasciicircum{}n)} if and only if
its order is prime to~$p$: the order~$0$ of an element of infinite order is not prime to a prime.
The implication from right to left holds in every commutative group
(\texttt{CommGroup.mem\_range\_powMonoidHom\_of\_coprime}); for the other, Krull's intersection
theorem gives $\bigcap_n p^n M = 0$ for every finitely generated $\ZZ_p$-module~$M$
(\texttt{PadicInt.eq\_zero\_of\_forall\_mem\_range\_nsmul\_pow}), and
\texttt{CommGroup.forall\_mem\_range\_powMonoidHom\_pow\_prod\_iff} treats the product with a finite
group.
\end{proof}

\begin{lemma}[The case $\ell = 2$]
\label{lem:ds-loc-two}
\lean{ChabautyColeman.HypJac.LocalStructure.mem_iInf_range_powMonoidHom_two_pow_iff}
\leanok
\uses{lem:ds-loc, def:hypjac}
Assume~(Loc) for a finite extension~$K_v$ of~$\QQ_2$. Then $\bigcap_{n \ge 0} 2^nJ(K_v)$ is the
subgroup of elements of odd order.
\end{lemma}

\begin{proof}
\leanok
Lemma~\ref{lem:ds-loc}\,(2) with $\ell = 2$.
\end{proof}

For odd~$\ell$, by contrast, $2$ is a unit of~$O_v$, and $\bigcap_n 2^nJ(K_v)$ contains a copy
of~$O_v^g$, so it is not torsion. The proof of Theorem~\ref{thm:ds-e1} uses
Lemma~\ref{lem:ds-loc-two} for $K_v = \QQ_2$.

\begin{definition}[Relaxed Selmer groups]
\label{def:ds-relaxed}
\lean{ChabautyColeman.relaxedSelmerGroup, ChabautyColeman.range_descentMap_le_relaxedSelmerGroup,
  ChabautyColeman.finite_relaxedSelmerGroup}
\leanok
\uses{def:ds-selmer, def:ds-mu, def:ds-sel2}
For a finite set $\Sigma \supseteq \Sigma_f$ of primes and a set~$V$ of places of~$K$, let
\[ S(\Sigma, V) = \{\xi \in A(\Sigma, 2) : N(\xi) = 1 \text{ and } \operatorname{res}_v(\xi)
   \in \mu_v(J(K_v)) \text{ for } v \in V\}\,, \]
with the local conditions of Definition~\ref{def:ds-sel2} at the places in~$V$ only. Then
$S(\Sigma, V)$ is finite, and $\mu(J(K)) \subseteq S(\Sigma, V)$.
\end{definition}

\begin{proof}
\leanok
\uses{thm:ds-image-selmer, thm:ds-norm, thm:ds-naturality}
Definition~\ref{def:ds-selmer} for the finiteness, and Theorems~\ref{thm:ds-image-selmer},
\ref{thm:ds-norm} and~\ref{thm:ds-naturality} for the inclusion.

\emph{Formalization.} \texttt{relaxedSelmerGroup} takes the places of~$V$ as a family of field
extensions of~$K$ (their completions), with the local conditions \texttt{HypJac.localCondition}
of \texttt{HypJac.selmerGroupTwo}; as these are taken at arbitrary field extensions, for $K = \QQ$
a condition at~$2$ can use Mathlib's field~$\QQ_2$ itself. The inclusion takes the hypothesis
that~$f$ has separable reduction outside~$\Sigma$ (\texttt{Polynomial.HasSeparableReduction}), as
Theorem~\ref{thm:ds-image-selmer} does; \texttt{finite\_relaxedSelmerGroup} is the finiteness for
finite~$\Sigma$, with the instances on class groups and unit groups of
\texttt{AdjoinRoot.finite\_selmerGroup}.
\end{proof}

For $K = \QQ$, $2 \in \Sigma$ and $V = (\Sigma \cup \{\infty\}) \setminus \{2\}$, $S(\Sigma, V)$ is
the group~$S$ of Remark~4.2 of [Sto17], which leaves out the $2$-adic condition, with~$\Sigma$ in
place of the set of bad primes for $2$-descent used there; $S(\Sigma, \emptyset) = A(\Sigma, 2) \cap
\ker N$ is coarser, and $\Sel_2 = S(\Sigma, V)$ when~$V$ is the set of all places, by
Theorem~\ref{thm:ds-sel2}\,(2). Each of them contains~$\Sel_2$, and so it can replace the
$2$-Selmer group in the algorithm of Section~4 of [Sto17], by Remark~4.2 there.

\begin{remark}[Local conditions are needed]
\label{rem:ds-dimension}
Let $K = \QQ$ and~$f$ irreducible, and let $(r_1, r_2)$ be the signature of the field~$L$, so that
$r_1 + 2r_2 = 2g + 1$. The unramified classes include $O_L^\times/(O_L^\times)^2$, of dimension
$r_1 + r_2$ (the roots of unity of~$L$ are~$\pm 1$, as $r_1 \ge 1$), on which the norm is onto
$\{\pm 1\}$, as $N(-1) = -1$; so the unramified classes of norm~$1$ have dimension at least
$r_1 + r_2 - 1$, whatever the class group of~$L$. At the real place,
$J(\mathbb R) \cong (\mathbb R/\ZZ)^g \times (\ZZ/2)^k$ and
$\#J(\mathbb R)[2] = 2^{r_1 + r_2 - 1}$ (Theorem~\ref{thm:ds-two-torsion} over~$\mathbb R$), so
$k = r_1 + r_2 - 1 - g = (r_1 - 1)/2$ is the dimension of $\mu_\infty(J(\mathbb R))$
(Theorem~\ref{thm:ds-kernel} over~$\mathbb R$), inside the space of dimension $r_1 - 1$ of sign
vectors of norm~$1$. So the classes that are unramified everywhere, of norm~$1$ and satisfy the
real condition form a space of dimension at least
$r_1 + r_2 - 1 - (r_1 - 1)/2 = (r_1 - 1)/2 + r_2 = g$, and so does every group defined by these
conditions with ramification allowed at some primes, such as $S(\Sigma, \{\infty\})$, which
contains that space. As $J(\QQ)[2] = 0$, the rank bound of Theorem~\ref{thm:ds-rank-bound} from
such a group is never below~$g$, and every example of rank less than~$g$ with irreducible~$f$
needs a local condition at~$2$ or at an odd bad prime, hence the local image size of
Lemma~\ref{lem:ds-loc}\,(1) and~(Loc).
\end{remark}

\begin{theorem}[The rank bound]
\label{thm:ds-rank-bound}
\lean{ChabautyColeman.HypJac.pow_rank_le_card_selmerGroupTwo,
  ChabautyColeman.pow_rank_le_card_relaxedSelmerGroup}
\leanok
\uses{def:ds-cited, def:ds-sel2, def:ds-relaxed}
Assume~(MW). Then $2^{\operatorname{rk} J(K)}\,\#J(K)[2] \le \#S$ for every finite subgroup
$S \subseteq L^\square$ containing~$\mu(J(K))$, in particular for $S = \Sel_2$ and for the
groups~$S(\Sigma, V)$ of Definition~\ref{def:ds-relaxed}.
\end{theorem}

\begin{proof}
\leanok
\uses{thm:ds-rank, thm:ds-sel2, def:ds-relaxed}
Theorem~\ref{thm:ds-rank}, with Theorem~\ref{thm:ds-sel2} and Definition~\ref{def:ds-relaxed} for
the inclusions and the finiteness.

\emph{Formalization.} Both statements take the instance \texttt{IsMulFG (HypJac K f)}, which the
class \texttt{HypJac.MordellWeil} provides over a number field, and the instances of the
finiteness statements; they are \texttt{HypJac.pow\_rank\_le\_card\_of\_range\_descentMap\_le} of
Theorem~\ref{thm:ds-rank} applied to the inclusions.
\end{proof}

\section{Examples}
\label{sec:ds-examples}

\begin{remark}[E1: genus~$2$, $f$ split]
\label{rem:ds-e1}
The first example is the curve
\[ C \colon y^2 = f(x) = x(x^2 - 1)(x^2 - 9) = x^5 - 10x^3 + 9x\,, \]
with the roots $\alpha = -3, -1, 0, 1, 3$ and $\operatorname{disc} f = 2^{16} \cdot 3^6$, so
$\Sigma_f = \{2, 3\}$. No smaller set is possible: among five integers, two are congruent
modulo~$2$ and two modulo~$3$. Here $L = \QQ^5$, and $A(\Sigma_f, 2) \cap \ker N$ has dimension
$4(\#\Sigma_f + 1) = 12$ over~$\FF_2$, whereas $\mu(J(\QQ)[2])$ has dimension~$4$. A local
condition at one place cuts out at most $4$ dimensions (an odd prime) or~$6$ (the prime~$2$), so
no single place suffices; the set $V = \{2, 3\}$ does: $S(\Sigma_f, V) = \mu(J(\QQ)[2])$
(Proposition~\ref{prop:ds-e1-sel}). Theorem~\ref{thm:ds-e1} shows $J(\QQ) = J(\QQ)[2] \cong
(\ZZ/2)^4$ and $C(\QQ) = \{\infty, (0, 0), (\pm 1, 0), (\pm 3, 0)\}$. Its proof needs neither~(MW)
nor any $2$-adic analysis on~$C$: the descent and~(Loc) for~$\QQ_2$ and~$\QQ_3$ give
$J(\QQ) = J(\QQ)[2] + T$ with~$T$ of odd order, and reduction at the good prime~$5$ kills~$T$.
\end{remark}

\begin{setting}
\label{set:ds-e1}
\lean{ChabautyColeman.DescentExample1.f, ChabautyColeman.DescentExample1.F,
  ChabautyColeman.DescentExample1.S}
\leanok
Let $f = x^5 - 10x^3 + 9x \in \QQ[x]$, with the roots
$\alpha_1, \dots, \alpha_5 = -3, -1, 0, 1, 3$, and $L = \QQ[T]/(f) \cong \QQ^5$,
$\xi \mapsto (\xi(\alpha_i))_i$. Let $\Sigma = \{2, 3\}$, let~$V$ consist of~$\QQ_2$ and~$\QQ_3$,
and $S = S(\Sigma, V)$ (Definition~\ref{def:ds-relaxed}). Let $\bar f = x^5 - x \in \FF_5[x]$ be
the reduction of~$f$ modulo~$5$.
\end{setting}

\begin{remark}[Formalization]
In the module \texttt{HyperellipticCurves.Descent.Example1}: \texttt{f} over~$\QQ$ and the
integral model \texttt{F} for the reduction modulo~$5$ (\texttt{DescentExample1.map\_F}), the
roots as a vector of integers, and~$\Sigma$ as~\texttt{S}, the primes $(2)$, $(3)$ of~$\ZZ$
(\texttt{Rat.HeightOneSpectrum.primesEquiv}). The family~$V$ consists of the $p$-adic fields of
Mathlib for $p = 2, 3$. $f$ has separable reduction outside~$\Sigma$
(\texttt{DescentExample1.hasSeparableReduction\_of\_notMem}), as
$\operatorname{disc} F = 2^{16} \cdot 3^6$ (\texttt{DescentExample1.discr\_F}, from Mathlib's
\texttt{Polynomial.resultant\_deriv}: the product of the values $144, -16, 9, -16, 144$ of~$F'$ at
the roots).
\end{remark}

\begin{lemma}[Square classes of local fields]
\label{lem:ds-padic-squares}
\lean{HenselianLocalRing.isSquare_iff_isSquare_residue, IsDiscreteValuationRing.modPowTwoEquiv,
  PadicInt.isSquare_iff_toZModPow_three_eq_one, Padic.modPowTwoEquivZModEight,
  Padic.modPowTwoEquiv_mk_of_eq_pow_mul}
\leanok
\uses{thm:hensel}
\begin{enumerate}
  \item Let~$O$ be a Henselian local ring in which~$2$ is a unit, with residue field~$k$. A unit
    of~$O$ is a square if and only if its image in~$k$ is a square.
  \item Let~$O$ moreover be a discrete valuation ring, with fraction field~$F$ and a
    uniformizer~$\pi$. Then $\pi^n u \mapsto (n \bmod 2, \bar u)$, for $u \in O^\times$, induces
    an isomorphism $F^\times/(F^\times)^2 \cong \ZZ/2 \times k^\times/(k^\times)^2$; for
    finite~$k$, both have order~$4$. This covers the finite extensions of~$\QQ_p$ for odd~$p$,
    among them~$\QQ_3$.
  \item A unit~$u$ of~$\ZZ_2$ is a square if and only if $u \equiv 1 \bmod 8$, and
    $2^n u \mapsto (n \bmod 2, u \bmod 8)$ induces an isomorphism
    $\QQ_2^\times/(\QQ_2^\times)^2 \cong \ZZ/2 \times (\ZZ/8)^\times$, of order~$8$.
\end{enumerate}
\end{lemma}

\begin{proof}
\leanok
(1) The image of a square is a square. Conversely, if $\bar u = \bar a^2$, Hensel's lemma applies
to $X^2 - u$ at~$a$, whose derivative~$2a$ is a unit. (2) $F^\times = \pi^{\ZZ} \times O^\times$,
and $\pi^n u$ is a square if and only if~$n$ is even and~$u$ is a square, which is~(1). (3) The
square of an odd integer is~$\equiv 1 \bmod 8$. Conversely, Hensel's lemma in the form for
$\lvert g(a)\rvert < \lvert g'(a)\rvert^2$ applies to $g = X^2 - u$ at $a = 1$, as
$\lvert 1 - u\rvert_2 \le 2^{-3} < \lvert 2\rvert_2^2$. The isomorphism follows as in~(2).

\emph{Formalization.} (1) is proved from the defining property of \texttt{HenselianLocalRing}; a
local ring that is adically complete at its maximal ideal, such as~$\ZZ_p$, is Henselian
(\texttt{IsAdicComplete.henselianLocalRing}). (2) comes from
$F^\times/(F^\times)^n \cong \ZZ/n \times O^\times/(O^\times)^n$ for every~$n$
(\texttt{IsDiscreteValuationRing.modPowEquiv}, from \texttt{IsDiscreteValuationRing.unitsEquivProd})
and $O^\times/(O^\times)^2 \cong k^\times/(k^\times)^2$ by~(1)
(\texttt{HenselianLocalRing.modPowTwoEquiv}), stated with the classes \texttt{Units.ModPow}; the
order~$4$ for finite~$k$ is not formalized, the examples compute with the isomorphism itself.
For~$\QQ_p$, $p$ odd, the isomorphism is \texttt{Padic.modPowTwoEquiv}, with the uniformizer~$p$
and the residue field $\ZZ/p$ (\texttt{PadicInt.residueField}); its value at~$p^n u$ for a unit~$u$
of~$\ZZ_p$ is \texttt{Padic.modPowTwoEquiv\_mk\_of\_eq\_pow\_mul}. (3) comes from Mathlib's
\texttt{hensels\_lemma}; the values of the isomorphisms on classes $p^n m$ with $m \in \ZZ$ are
\texttt{Padic.modPowTwoEquivZModEight\_mk\_of\_val\_eq} and
\texttt{Padic.modPowTwoEquiv\_mk\_of\_val\_eq}. Mathlib has none of these statements.
\end{proof}

\begin{lemma}[$A(\Sigma, 2)$ for split~$f$ over~$\QQ$]
\label{lem:ds-e1-selmer}
\lean{Rat.selmerGroup_eq_closure, AdjoinRoot.mem_selmerGroup_iff_of_splits}
\leanok
\uses{def:ds-selmer}
\begin{enumerate}
  \item For a finite set~$\Sigma$ of primes, $\QQ(\Sigma, n)$ is generated by the classes of~$-1$
    and of the primes in~$\Sigma$.
  \item If $f \in \QQ[x]$ is the product of distinct factors $x - \alpha_i$, a class
    $\xi \in L^\times/(L^\times)^n$ lies in $A(\Sigma, n)$ if and only if every
    component~$\xi(\alpha_i)$ lies in~$\QQ(\Sigma, n)$.
\end{enumerate}
\end{lemma}

\begin{proof}
\leanok
\uses{lem:ds-norm-selmer}
(1) Unique factorization in~$\ZZ$, with $\ZZ^\times = \{\pm 1\}$. (2) The field factors of~$L$ are
the $\QQ[x]/(x - \alpha_i) \cong \QQ$, with integral closure~$\ZZ$, and a class of~$\QQ$ lies
in~$\QQ(\Sigma, n)$ if and only if its image in such a factor lies in its Selmer group: one
direction is restriction, the other the norm (Lemma~\ref{lem:ds-norm-selmer}), here of degree~$1$.

\emph{Formalization.} (1) is stated for a finite set of primes of~$\ZZ$, with Mathlib's
\texttt{Rat.HeightOneSpectrum.natGenerator} (the prime number of a prime of~$\ZZ$), as an equality
of subgroups; the proof adjoins the primes one at a time
(\texttt{IsDedekindDomain.selmerGroup\_insert}) to $\ZZ^\times/(\ZZ^\times)^n$, which is all
of~$\QQ(\emptyset, n)$ as~$\ZZ$ has trivial class group. (2) is stated for a field~$K$ with a
Dedekind domain~$R$ (the field factors of $K[x]/(f)$ are copies of~$K$, the integral closure
is~$R$, \texttt{IsDedekindDomain.map\_algebraMap\_mem\_selmerGroupAbove\_iff} for degree~$1$), and
$f$ with at least $\deg f$ distinct roots. The components~$\xi(\alpha_i)$ are those of the
isomorphism \texttt{AdjoinRoot.modPowEquivPiOfRoots} from $L^\times/(L^\times)^n$ to the product
of copies of $K^\times/(K^\times)^n$, given by evaluation at the roots
(\texttt{AdjoinRoot.equivPiOfRoots}, Lagrange interpolation); it commutes with base change
(\texttt{AdjoinRoot.modPowEquivPiOfRoots\_map}).
\end{proof}

\begin{lemma}[The local images of E1]
\label{lem:ds-e1-local}
\lean{ChabautyColeman.DescentExample1.range_descentMap_three,
  ChabautyColeman.DescentExample1.range_descentMap_two}
\leanok
\uses{lem:ds-loc, thm:ds-two-torsion, thm:ds-kernel, thm:ds-formula,
  cor:ds-torsion-values, lem:ds-padic-squares}
Assume~(Loc) for~$\QQ_2$ and~$\QQ_3$. Then $\mu_3(J(\QQ_3))$ is the group of order~$16$ generated
by the images of the five points $(\alpha_j, 0)$, and $\mu_2(J(\QQ_2))$ is the group of order~$64$
generated by these and the images of the points of~$C(\QQ_2)$ with $x = 4$ and $x = 1/5$.
\end{lemma}

\begin{proof}
\leanok
The group $\mu_v(J(\QQ_v)) \cong J(\QQ_v)/2J(\QQ_v)$ (Theorem~\ref{thm:ds-kernel} over~$\QQ_v$) has
order $\#J(\QQ_v)[2] = 16$ for $v = 3$ and $2^2 \cdot 16 = 64$ for $v = 2$, by
Lemma~\ref{lem:ds-loc}\,(1) and Theorem~\ref{thm:ds-two-torsion}, as~$f$ has five linear factors
over~$\QQ_v$. The listed classes lie in it, the classes of $4 - \theta$ and $1/5 - \theta$ by
Theorem~\ref{thm:ds-formula}: the two points over~$\QQ_2$ exist by
Lemma~\ref{lem:ds-padic-squares}\,(3), as $f(4) = 2^2 \cdot 105$ and
$f(1/5) = 2^8 \cdot 21/5^5$ with $105 \equiv 21 \cdot 5^5 \equiv 1 \bmod 8$. The isomorphisms of
Lemma~\ref{lem:ds-padic-squares}\,(2) for~$\QQ_3$ and~(3) for~$\QQ_2$, applied to the five
components of the values (Corollary~\ref{cor:ds-torsion-values} and
Theorem~\ref{thm:ds-formula}), separate~$16$, resp.~$64$, products of the listed classes.

\emph{Formalization.} The classes of $4 - \theta$ and $1/5 - \theta$ enter as units of
$\QQ_2[T]/(f)$ with these values; they lie in the image by
\texttt{HypJac.descentMap\_mumfordClass\_X\_sub\_C}, the value at a point $(x_0, y_0)$ with
$f(x_0) \ne 0$. The components of the values $\mu((\alpha_j, 0))$ at the roots~$\alpha_i$ are
$\operatorname{lc}(f)(\alpha_j - \alpha_i)$ for $i \ne j$ and $\operatorname{lc}(f) f'(\alpha_j)$
for $i = j$ (Corollary~\ref{cor:ds-torsion-values}), written as products of~$-1$, $2$, $3$ with a
table of exponents that \texttt{decide} checks against the integers. The order of the image is
\texttt{HypJac.card\_range\_descentMap\_eq\_card\_quotient} with Lemma~\ref{lem:ds-loc}; the lower
bound for the subgroup generated by the listed classes is
\texttt{Subgroup.card\_closure\_range\_eq\_two\_pow} (elements of order dividing~$2$ of which no
nonempty product is trivial generate a group of order~$2^k$), applied to the images in root
coordinates, where \texttt{decide} checks the products. Four of the five points $(\alpha_j, 0)$
suffice: the five values have product~$1$.
\end{proof}

\begin{proposition}[The relaxed Selmer group of~E1]
\label{prop:ds-e1-sel}
\lean{ChabautyColeman.DescentExample1.relaxedSelmerGroup_eq}
\leanok
\uses{def:ds-relaxed, lem:ds-e1-selmer, lem:ds-e1-local}
Assume~(Loc) for~$\QQ_2$ and~$\QQ_3$. Then $S(\Sigma, V) = \mu(J(\QQ)[2])$.
\end{proposition}

\begin{proof}
\leanok
\uses{cor:ds-torsion-values, thm:ds-naturality, lem:ds-padic-squares}
The inclusion~$\supseteq$ is Definition~\ref{def:ds-relaxed}, as~$f$ has separable reduction
outside~$\Sigma$. Let $\xi \in S(\Sigma, V)$. By Lemma~\ref{lem:ds-e1-selmer}, the components of
the elements of~$A(\Sigma, 2)$ lie in the group generated by the classes of~$-1$, $2$ and~$3$,
which stay independent in~$\QQ_2^\times/(\QQ_2^\times)^2$ (Lemma~\ref{lem:ds-padic-squares}\,(3));
so the map from~$A(\Sigma, 2)$ to $L_2^\times/(L_2^\times)^2$, $L_2 = \QQ_2[T]/(f)$, is injective.
The classes~$P$ and~$Q$ of~$A(\Sigma, 2)$ with components $(-1, -3, 1, 3, 1)$ and
$(-3, -2, -3, 3, -6)$ map to the classes of $4 - \theta$ and $1/5 - \theta$, as the components
agree in~$\QQ_2^\times/(\QQ_2^\times)^2$. By Lemma~\ref{lem:ds-e1-local} for~$\QQ_2$, as~$\mu$
commutes with the map to~$\QQ_2$ (Theorem~\ref{thm:ds-naturality}), $\xi$ therefore lies in the
group generated by~$P$, $Q$ and the values $\mu((\alpha_j, 0))$
(Corollary~\ref{cor:ds-torsion-values}), which lie in~$A(\Sigma, 2)$: $\xi = wz$ with
$w \in \mu(J(\QQ)[2])$ and $z \in \{1, P, Q, PQ\}$. Both~$\xi$ and~$w$ satisfy the local condition
at~$\QQ_3$, hence so does~$z$. By Lemma~\ref{lem:ds-e1-local} for~$\QQ_3$, the image of~$z$ is
then a product of the images of the values $\mu((\alpha_j, 0))$, which the isomorphism of
Lemma~\ref{lem:ds-padic-squares}\,(2), applied to the five components, excludes for $z \ne 1$.
So $\xi = w$.

\emph{Formalization.} The classes of~$A(\Sigma, 2)$ and their images over~$\QQ_2$ and~$\QQ_3$ are
compared in root coordinates (\texttt{AdjoinRoot.modPowEquivPiOfRoots} and its naturality
\texttt{AdjoinRoot.modPowEquivPiOfRoots\_map}), every component a product of the classes of~$-1$,
$2$, $3$ given by a table of exponents, and the last step is a check by \texttt{decide}. The
elements of a subgroup generated by elements of order dividing~$2$ are the products over the finite
sets of generators (\texttt{Subgroup.mem\_closure\_range\_iff\_of\_sq\_eq\_one}). The norm
condition is not used.
\end{proof}

\begin{definition}[A cited input: reduction]
\label{def:ds-red}
\lean{ChabautyColeman.HypJac.TorsionReduction}
\leanok
\begin{description}
  \item[\textmd{(Red)}] Let $F \in \ZZ[x]$ be monic of odd degree and squarefree, and~$p \ge 3$
    a prime not dividing $\operatorname{disc} F$. Then reduction modulo~$p$ maps the torsion
    subgroup of~$J(\QQ)$ injectively into~$\tilde J(\FF_p)$, the Jacobian of $y^2 = \bar F(x)$
    (N.~M.~Katz, \emph{Galois properties of torsion points on abelian varieties}, Invent.\
    Math.~62 (1981), 481--502, Appendix: the kernel of reduction is torsion-free, as
    $e = 1 < p - 1$; for torsion of order prime to~$p$ also M.~Hindry and J.~H.~Silverman,
    \emph{loc.\ cit.}, Theorem~C.1.4).
\end{description}
\end{definition}

\begin{remark}[Formalization]
\texttt{HypJac.TorsionReduction p F}, for $F \in \ZZ[x]$, is a structure with values in
\texttt{Prop}, like~(Loc): there is a homomorphism from the Jacobian over~$\QQ$ of the image
of~$F$ to the Jacobian over~$\ZZ/p$ of its reduction that is injective on the elements of finite
order, in the form that an element of finite order in its kernel is trivial. The reduction map
is not defined; the example uses only such a homomorphism, through
\texttt{HypJac.TorsionReduction.pow\_eq\_one\_of\_isOfFinOrder}: an exponent of the group over
$\ZZ/p$ kills every element of finite order of~$J(\QQ)$.
\end{remark}

\begin{lemma}[$\tilde J(\FF_5)$]
\label{lem:ds-e1-five}
\lean{ChabautyColeman.DescentExample1.exponent_hypJac_zmod_five,
  ChabautyColeman.DescentExample1.sq_eq_one_of_eq_X_pow_five_sub_X}
\leanok
\uses{def:hypjac}
The group $\tilde J(\FF_5)$ of $y^2 = \bar f(x) = x^5 - x$ has exponent~$2$; so it is
$\tilde J(\FF_5)[2]$, of order~$16$.
\end{lemma}

\begin{proof}
\leanok
\uses{thm:ds-two-torsion, lem:mumford-reduce}
Every class has a reduced Mumford representative $(u, v)$: $u$ monic, $\deg v < \deg u \le 2$ and
$u \mid v^2 - \bar f$ (Lemma~\ref{lem:mumford-reduce}). For $u = x - a$, $v^2 = \bar f(a) = 0$.
For $u = x^2 + bx + c$ and $v = dx + e$, the remainder of $v^2 - \bar f$ modulo~$u$ has two
coefficients, polynomials in $b, c, d, e$, which vanish over~$\FF_5$ only for $d = e = 0$
(\texttt{decide} over the~$5^4$ tuples). So $v = 0$, and the class of $(u, 0)$ has order
dividing~$2$. The order~$16$ is Theorem~\ref{thm:ds-two-torsion}.

\emph{Formalization.} The statements are for any polynomial over $\ZZ/5$ equal to $x^5 - x$, so
that they apply to the reduction of the integral model (\texttt{DescentExample1.map\_F\_zmod\_five},
by \texttt{reduce\_mod\_char}). The form that Theorem~\ref{thm:ds-e1} consumes is
\texttt{DescentExample1.sq\_eq\_one\_of\_eq\_X\_pow\_five\_sub\_X}; the order~$16$ is not stated.
\end{proof}

\begin{theorem}[E1]
\label{thm:ds-e1}
\lean{ChabautyColeman.DescentExample1.ker_powMonoidHom_two_eq_top,
  ChabautyColeman.DescentExample1.card_hypJac, ChabautyColeman.DescentExample1.ratPoints_eq}
\leanok
\uses{def:ds-cited, def:ds-red, def:hypjac}
Assume~(Loc) for~$\QQ_2$ and~$\QQ_3$ and~(Red) for $p = 5$. Then
$J(\QQ) = J(\QQ)[2] \cong (\ZZ/2)^4$, and $C(\QQ) = \{\infty, (0, 0), (\pm 1, 0), (\pm 3, 0)\}$.
\end{theorem}

\begin{proof}
\leanok
\uses{prop:ds-e1-sel, thm:ds-kernel, def:ds-relaxed, lem:ds-loc, lem:ds-loc-two,
  thm:basechange-inj, lem:ds-e1-five, thm:ds-two-torsion, lem:mumford-exists}
Write $J = J(\QQ)$ and $\Gamma = J[2]$, a group of order~$16$ by Theorem~\ref{thm:ds-two-torsion}.
By Definition~\ref{def:ds-relaxed} and Proposition~\ref{prop:ds-e1-sel},
$\mu(J) \subseteq \mu(\Gamma)$, and as $\ker\mu = 2J$ (Theorem~\ref{thm:ds-kernel}),
$J = \Gamma + 2J$. Multiplying by~$2^n$ gives $2^nJ = 2^n\Gamma + 2^{n+1}J \subseteq \Gamma +
2^{n+1}J$, so $J = \Gamma + 2^nJ$ for every~$n$, by induction. Let $a \in J$, and for each~$n$
choose $\gamma_n \in \Gamma$ with $a - \gamma_n \in 2^nJ$. As~$\Gamma$ is finite, some
$\gamma \in \Gamma$ equals~$\gamma_n$ for infinitely many~$n$; as the subgroups~$2^nJ$ decrease,
$t = a - \gamma$ lies in $\bigcap_n 2^nJ$. Every element of $\bigcap_n 2^nJ$ has odd order: its
image in~$J(\QQ_2)$ does, by Lemma~\ref{lem:ds-loc-two}, and $J \to J(\QQ_2)$ is injective
(Theorem~\ref{thm:basechange-inj}). By~(Red) and Lemma~\ref{lem:ds-e1-five}, the image of~$t$
in~$\tilde J(\FF_5)$ has odd order and order dividing~$2$, so it is~$0$, and $t = 0$. So
$J = \Gamma$. Finally let $P = (x_0, y_0) \in C(\QQ)$. As $[P - \infty]$ has order dividing~$2$,
$[P - \infty] = [(x_0, -y_0) - \infty]$; the reduced Mumford representatives $(x - x_0, y_0)$
and $(x - x_0, -y_0)$ of these classes agree (Lemma~\ref{lem:mumford-exists}), so $y_0 = 0$ and
$f(x_0) = 0$.

\emph{Formalization.} In the module \texttt{HyperellipticCurves.Descent.Example1}. $C(\QQ)$ is
stated for the affine points, $\{(x, y) \in \QQ^2 : y^2 = f(x)\} = \{(\alpha_i, 0)\}$. The
hypotheses are \texttt{HypJac.LocalStructure} for~$\QQ_2$ and~$\QQ_3$ and
\texttt{HypJac.TorsionReduction 5}; \texttt{HypJac.MordellWeil} does not occur.
$J(\QQ) = J(\QQ)[2] \cong (\ZZ/2)^4$ is stated as \texttt{ker\_powMonoidHom\_two\_eq\_top}, the
kernel of the squaring map is everything, with \texttt{card\_hypJac}, $\#J(\QQ) = 16$. The step
from $\Gamma + 2J$ to $\Gamma + 2^nJ$ is \texttt{Subgroup.sup\_range\_powMonoidHom\_pow\_eq\_top},
and the reduced representatives are compared by \texttt{mumfordClass\_inj}. The Lean concludes
through the torsion of~$J(\QQ)$: by the argument with the finite group~$\Gamma$
(\texttt{Subgroup.coprime\_orderOf\_mk\_of\_mem\_iInf\_sup\_range\_powMonoidHom\_pow}),
$J(\QQ)/\Gamma$ is a torsion group, hence so is~$J(\QQ)$ (Mathlib's
\texttt{IsMulTorsion.extension\_closed}), and by~(Red) and Lemma~\ref{lem:ds-e1-five} every element
of~$J(\QQ)$ has order dividing~$2$.
\end{proof}

\begin{remark}[E2: rank~$1$]
\label{rem:ds-e2}
The second example is the curve
\[ C \colon y^2 = f(x) = (x + 4)(x + 3)(x^3 - x + 1) = x^5 + 7x^4 + 11x^3 - 6x^2 - 5x + 12\,, \]
with $\operatorname{disc} f = -23^3 \cdot 59^2$, so $\Sigma_f = \{23, 59\}$. Here
$L = \QQ \times \QQ \times F$, where $F = \QQ[t]/(t^3 - t + 1)$ is the cubic field of
discriminant~$-23$, of class number~$1$ and unit rank~$1$. The places $V = \{23, 59\}$ suffice for
the descent: the prime~$2$ is not needed, and the components in~$\QQ_v$ of the local conditions
suffice, apart from one class at~$23$ (Lemma~\ref{lem:ds-e2-local}).

The group $J(\QQ)[2] \cong (\ZZ/2)^2$ (Theorem~\ref{thm:ds-two-torsion}) is the torsion subgroup
of~$J(\QQ)$ (by a computation with Magma), and the class of $(-11/9, 1020/729) - \infty$ has
infinite order (Lemma~\ref{lem:ds-e2-points}). Under~(MW) and~(Loc) for~$\QQ_{23}$ and~$\QQ_{59}$,
$\operatorname{rk} J(\QQ) = 1$, and $\Sel_2 = S(\Sigma_f, V)$ has dimension~$3$, as it lies between
$\mu(J(\QQ))$ and $S(\Sigma_f, V)$ (Theorem~\ref{thm:ds-e2}). The known rational points are
$\infty$, $(-4, 0)$, $(-3, 0)$ and $(-11/9, \pm 1020/729)$. As $\operatorname{rk} J(\QQ) = 1 < g$,
Theorem~\ref{thm:main} with $r = 1$ gives $\#C(\QQ) \le \#\bar C(\FF_p) + 2$ for the primes
$p \ge 5$ of good reduction; this bound is at least~$8$ at every such prime (it is~$8$ for $p = 5$
and $p = 7$, and for $p > 19$ the Weil bound gives more). So it does not determine~$C(\QQ)$; that
needs further methods, such as the Selmer group Chabauty method of [Sto17], and is not formalized
here.
\end{remark}

\begin{setting}
\label{set:ds-e2}
\lean{ChabautyColeman.DescentExample2.g, ChabautyColeman.DescentExample2.f,
  ChabautyColeman.DescentExample2.S, ChabautyColeman.DescentExample2.P}
\leanok
Let $g = x^3 - x + 1$ and $f = (x + 4)(x + 3)g \in \QQ[x]$, let $F = \QQ[t]/(g)$, with ring of
integers~$O_F$, and $L = \QQ[T]/(f) \cong \QQ \times \QQ \times F$,
$\xi \mapsto (\xi(-4), \xi(-3), \xi(t))$; we write classes in $L^\times/(L^\times)^2$ as triples
of their components. Let $\Sigma = \{23, 59\}$, let~$V$ consist of~$\QQ_{23}$ and~$\QQ_{59}$, and
$S = S(\Sigma, V)$ (Definition~\ref{def:ds-relaxed}). Let $T_1, T_2, P \in J(\QQ)$ be the classes
of $(-4, 0) - \infty$, $(-3, 0) - \infty$ and $(-11/9, 1020/729) - \infty$.
\end{setting}

\begin{remark}[Formalization]
In the modules \texttt{HyperellipticCurves.Descent.Example2.Basic} (Setting~\ref{set:ds-e2} to
Lemma~\ref{lem:ds-e2-padic}\,(1), and Lemma~\ref{lem:ds-e2-points}) and
\texttt{HyperellipticCurves.Descent.Example2.Rank} (Lemma~\ref{lem:ds-e2-local},
Proposition~\ref{prop:ds-e2-sel}, Theorem~\ref{thm:ds-e2}), as for~E1. $f$ has separable reduction
outside~$\Sigma$, also at~$2$, where \texttt{Polynomial.HasSeparableReduction} asks nothing special:
$\operatorname{disc} f = -23^3 \cdot 59^2$, from
$\operatorname{disc}(pq) = \operatorname{disc}(p) \operatorname{disc}(q) \operatorname{Res}(p, q)^2$
with $\operatorname{disc} g = -23$, $g(-4) = -59$ and $g(-3) = -23$. The classes are
\texttt{mumfordClass} values, $P$ that of $(x + 11/9, 1020/729)$. The factors $x + 4$, $x + 3$, $g$
of~$f$ are indexed by $\{0, 1, 2\}$ (\texttt{DescentExample2.factorsEquiv}).
\end{remark}

\begin{lemma}[The cubic field]
\label{lem:ds-e2-field}
\lean{ChabautyColeman.DescentExample2.subsingleton_classGroup,
  ChabautyColeman.DescentExample2.card_modPow_two,
  ChabautyColeman.DescentExample2.range_map_eq_closure,
  ChabautyColeman.DescentExample2.selmerGroupAbove_eq_closure}
\leanok
\uses{def:ds-selmer}
\begin{enumerate}
  \item $O_F$ has trivial class group, and $O_F^\times/(O_F^\times)^2$ has order~$4$ and is
    generated by the classes of~$-1$ and~$t$, a unit with $t^{-1} = 1 - t^2$.
  \item The elements $t + 3$, $2t - 3$, $t + 4$, $5t + 6$ and $t^2 + 4t - 3$ of~$O_F$, of norms
    $23$, $-23$, $59$, $-59$ and~$-59$, generate prime ideals, and these are the primes of~$O_F$
    above~$23$ and~$59$:
    \[ 23\,t^3 = -(t + 3)(2t - 3)^2\,, \qquad 59 = t^6\,(t + 4)(5t + 6)(t^2 + 4t - 3)\,. \]
  \item $F(\Sigma, 2)$ is generated by the classes of~$-1$, $t$ and the five elements of~(2).
\end{enumerate}
\end{lemma}

\begin{proof}
\leanok
\uses{lem:ds-e2-padic, lem:ds-padic-squares, lem:ds-resultant}
(1) The power basis of~$t$ has discriminant $\operatorname{disc} g = -23$, so the field
discriminant divides~$-23$ and is negative; the Minkowski bound gives class number~$1$, and
Dirichlet's unit theorem gives unit rank~$1$, the signature being~$(1, 1)$. As the roots of unity
are~$\pm 1$, $O_F^\times/(O_F^\times)^2$ has order~$2^{1 + 1}$. The classes of~$-1$, $t$ and~$-t$
are nontrivial: under $t \mapsto r_1$, resp.~$r_2$, the roots of~$g$ in~$\QQ_{59}$ with
$r_1 \equiv -4$ and $r_2 \equiv -13 \bmod 59$ (Lemma~\ref{lem:ds-e2-padic}), the images $-1$,
$r_1$ and~$-r_2$ are units that are not squares modulo~$59$
(Lemma~\ref{lem:ds-padic-squares}\,(2)). (2) The norm of $at + b$ is $b^3 - a^2b - a^3$
(Lemma~\ref{lem:ds-resultant}), and an element of norm~$\pm p$, for a prime~$p$, generates a prime
ideal of absolute norm~$p$. The norm of~$t^2 + 4t - 3$ follows from the identity for~$59$ by
multiplicativity, as $N(t) = -1$. A prime of~$O_F$ above~$23$ contains $23t^3$, hence $t + 3$ or
$2t - 3$, and so equals the prime ideal generated by it; likewise for~$59$. (3) As the class group
is trivial, $F(\emptyset, 2)$ is the image of the units, given by~(1), and adding the primes
of~(2) one at a time adds their generators.

\emph{Formalization.} $O_F$ is the integral closure of~$\ZZ$ in the field factor \texttt{AdjoinRoot
g} of~$L$. (1) uses general lemmas on number fields that are missing from Mathlib: for the class
number, the Minkowski bound for cubic fields
(\texttt{NumberField.RingOfIntegers.isPrincipalIdealRing\_of\_finrank\_eq\_three\_of\_abs\_discr\_le})
with \texttt{NumberField.isPrincipalIdealRing\_integralClosure\_iff}, and for the unit rank
\texttt{NumberField.Units.finrank\_eq\_one\_of\_finrank\_eq\_three\_of\_discr\_neg} with
\texttt{NumberField.finrank\_additive\_units\_integralClosure}; the order~$4$ is
\texttt{Units.card\_modPow\_two}. The classes of~$-1$, $t$, $-t$ are nonsquares by
\texttt{Padic.isSquare\_iff\_modPowTwoCoords\_eq\_zero} after the evaluations $t \mapsto r_1, r_2$
(\texttt{AdjoinRoot.lift}). (2) uses \texttt{NumberField.prime\_of\_abs\_norm\_eq} (an element of
norm~$\pm p$ is prime); the norms of linear elements come from
\texttt{AdjoinRoot.leadingCoeff\_mul\_norm\_of\_sub\_root}. (3) comes from
\texttt{IsDedekindDomain.selmerGroup\_image\_eq\_range\_sup\_closure} and
\texttt{IsDedekindDomain.selmerGroup\_range\_eq\_range\_sup\_closure}, the generalization of
Lemma~\ref{lem:ds-e1-selmer}\,(1): for trivial class group and a finite set of primes with given
generators, $K(\Sigma, n)$ is generated by the image of~$R^\times$ and the generators (by
\texttt{IsDedekindDomain.selmerGroup\_insert} and
\texttt{IsDedekindDomain.selmerGroup.fromUnit\_surjective}); \texttt{Rat.selmerGroup\_eq\_closure}
is proved through them. The consumed form of~(3),
\texttt{DescentExample2.selmerGroupAbove\_eq\_closure}, is for any polynomial equal to~$g$, so that
it applies to the field factor of~$f$ at the factor~$g$.
\end{proof}

\begin{lemma}[$A(\Sigma, 2)$ of~E2]
\label{lem:ds-e2-selmer}
\lean{ChabautyColeman.DescentExample2.selmerGroup_le_closure,
  ChabautyColeman.DescentExample2.card_selmerGroup_le}
\leanok
\uses{lem:ds-e1-selmer, lem:ds-e2-field}
A class $\xi \in L^\times/(L^\times)^2$ lies in $A(\Sigma, 2)$ if and only if
$\xi(-4), \xi(-3) \in \QQ(\Sigma, 2)$ and $\xi(t) \in F(\Sigma, 2)$. So $A(\Sigma, 2)$ is
generated by the $13$ classes with two trivial components and the third one equal to~$-1$, $23$
or~$59$ (at~$-4$ or~$-3$) or to one of the seven generators of Lemma~\ref{lem:ds-e2-field}\,(3)
(in~$F$), and $\#A(\Sigma, 2) \le 2^{13}$.
\end{lemma}

\begin{proof}
\leanok
\uses{def:ds-selmer}
The field factors of~$L$ are those of the three irreducible factors of~$f$, and $A(\Sigma, 2)$ is
defined componentwise (Definition~\ref{def:ds-selmer}). For the two linear factors,
Lemma~\ref{lem:ds-e1-selmer}\,(2) identifies the condition with
$\xi(-4), \xi(-3) \in \QQ(\Sigma, 2)$, generated by~$-1$, $23$, $59$
(Lemma~\ref{lem:ds-e1-selmer}\,(1)); for~$g$, Lemma~\ref{lem:ds-e2-field}\,(3).

\emph{Formalization.} \texttt{Polynomial.Factors.equivOfProd} identifies the factors of a product
of distinct monic irreducible polynomials with their index set, as
\texttt{Polynomial.Factors.equivOfRoots} does for linear factors. The $13$ classes are the images,
under the inverse of \texttt{AdjoinRoot.modPowEquivPiFactors}, of families with one nontrivial
component. A subgroup generated by~$n$ elements of order dividing~$2$ has at most~$2^n$ elements
(\texttt{Subgroup.card\_closure\_range\_le\_two\_pow}). Only the inclusion of $A(\Sigma, 2)$ in the
subgroup generated by the $13$ classes is formalized, which is what
Proposition~\ref{prop:ds-e2-sel} uses. The classes are given by data, the factor and the
coefficients of the nontrivial component, so that their coordinates in
Proposition~\ref{prop:ds-e2-sel} are computed by one lemma.
\end{proof}

\begin{lemma}[The local fields of E2]
\label{lem:ds-e2-padic}
\lean{PadicInt.exists_aeval_eq_zero_toZModPow_eq, Padic.not_isSquare_adjoinRoot_of_odd_valuation,
  Padic.modPowTwoEquiv_mk_of_toZModPow_eq, Padic.modPowTwoEquivZModTwo_mk_of_toZModPow_two_eq,
  ChabautyColeman.DescentExample2.exists_isRootData_fiftyNine,
  ChabautyColeman.DescentExample2.exists_isRootData_twentyThree}
\leanok
\uses{lem:ds-padic-squares}
\begin{enumerate}
  \item $g$ has a root $r \in \ZZ_{23}$ with $r \equiv 181 \bmod 23^2$, and roots
    $r_1, r_2, r_3 \in \ZZ_{59}$ with $r_1, r_2, r_3 \equiv 3182, 2406, 1374 \bmod 59^2$; so
    $r \equiv -3 \bmod 23$ and $r_1, r_2, r_3 \equiv -4, -13, 17 \bmod 59$.
  \item Over~$\QQ_{59}$, $f$ has the five distinct roots $-4, -3, r_1, r_2, r_3$.
    Over~$\QQ_{23}$, $f = (x + 4)(x + 3)(x - r)h$ with $h = x^2 + rx + r^2 - 1$ irreducible: its
    discriminant $4 - 3r^2 \equiv 5 \cdot 23 \bmod 23^2$ has odd valuation.
  \item Let~$p$ be odd and $x \in \ZZ_p$ with $x \equiv p^n m \bmod p^{n+1}$ for an integer~$m$
    prime to~$p$. Then~$x$ and $p^n m$ have the same class in $\QQ_p^\times/(\QQ_p^\times)^2$. In
    particular, for $Q \in \ZZ[x]$ and a root~$\rho$ of~(1) with $\rho \equiv a \bmod p^2$, the
    class of~$Q(\rho)$ is determined by $Q(a) \bmod p^2$ if $p^2 \nmid Q(a)$.
  \item Let $E = \QQ_{23}[x]/(h)$, with~$\theta_h$ the class of~$x$, and let
    $u, w \in \QQ_{23}$ with $u - wr/2 \in \ZZ_{23}^\times$ not a square modulo~$23$. Then
    $u + w\theta_h$ is not a square in~$E$. In particular, $-\theta_h$ is not a square, as
    $r/2 \equiv 10 \bmod 23$.
\end{enumerate}
\end{lemma}

\begin{proof}
\leanok
(1) Hensel's lemma, in the form with $\lvert \rho - a\rvert = \lvert g(a)\rvert$ for
$\lvert g'(a)\rvert = 1$, as $g(a) \equiv 0 \bmod p^2$ and $p \nmid g'(a)$ for these~$a$. (2) The
roots are distinct modulo~$59$ except $r_1 \equiv -4$, and $r_1 \ne -4$ as $g(-4) = -59$. A
root~$x_0$ of~$h$ in~$\QQ_{23}$ would make $4 - 3r^2 = (2x_0 + r)^2$ a square. (3) $x/(p^n m)$ is
a unit $\equiv 1 \bmod p$, hence a square (Lemma~\ref{lem:ds-padic-squares}\,(1)). (4) Let
$s = 2\theta_h + r$, so that $s^2 = D = 4 - 3r^2$ and
$u + w\theta_h = (u - wr/2) + (w/2)s$. If $u + w\theta_h = (\alpha + \beta s)^2$ with
$\alpha, \beta \in \QQ_{23}$, then $u - wr/2 = \alpha^2 + D\beta^2$. The two terms have
valuations of different parity, so~$\alpha^2$ is a unit and $D\beta^2 \in 23\ZZ_{23}$, and
$u - wr/2 \equiv \alpha^2 \bmod 23$.

\emph{Formalization.} (1) is Hensel's lemma with a congruence modulo~$p^k$
(\texttt{PadicInt.exists\_aeval\_eq\_zero\_toZModPow\_eq}, from Mathlib's
\texttt{hensels\_lemma}). (3) is the value of \texttt{Padic.modPowTwoEquiv} at $p^n u$ for a
unit~$u$ of~$\ZZ_p$ with a given residue (\texttt{Padic.modPowTwoEquiv\_mk\_of\_eq\_pow\_mul}, which
\texttt{Padic.modPowTwoEquiv\_mk\_of\_val\_eq} uses) and at $x \equiv p^n m \bmod p^k$, $n < k$
(\texttt{Padic.modPowTwoEquiv\_mk\_of\_toZModPow\_eq}). (4) is stated for a monic quadratic
over~$\QQ_p$, $p$ odd, whose discriminant has odd valuation
(\texttt{Padic.not\_isSquare\_adjoinRoot\_of\_odd\_valuation}); the irreducibility in~(2) comes from
\texttt{Polynomial.irreducible\_X\_sq\_add\_C\_mul\_X\_add\_C\_of\_not\_isSquare}. The example
uses~(3) in the form \texttt{Padic.modPowTwoEquivZModTwo\_mk\_of\_toZModPow\_two\_eq}: the class
of $c^2 y$, for $y \in \ZZ_p$ with residue $z \ne 0$ modulo~$p^2$, has the coordinates
\texttt{Padic.modPowTwoCoords p z} in~$(\ZZ/2)^2$ (the valuation modulo~$2$, and whether the unit
part is a nonsquare modulo~$p$), which \texttt{decide} computes. The roots enter only inside
proofs, through \texttt{DescentExample2.exists\_isRootData\_fiftyNine} and
\texttt{DescentExample2.exists\_isRootData\_twentyThree}; the lemmas on coordinates take a vector
of roots with their residues modulo~$p^2$ (\texttt{DescentExample2.IsRootData}), and no main
statement mentions them.
\end{proof}

\begin{lemma}[The local images of E2]
\label{lem:ds-e2-local}
\lean{ChabautyColeman.DescentExample2.range_descentMap_twentyThree,
  ChabautyColeman.DescentExample2.range_descentMap_fiftyNine}
\leanok
\uses{lem:ds-loc, thm:ds-two-torsion, thm:ds-kernel, thm:ds-formula, cor:ds-torsion-values,
  thm:ds-naturality, lem:ds-e2-padic}
Assume~(Loc) for~$\QQ_{23}$ and~$\QQ_{59}$. Then $\mu_{59}(J(\QQ_{59}))$ is the group of
order~$16$ generated by the images of~$T_1$, $T_2$ and of the points of~$C(\QQ_{59})$ with
$x = -1/2$ and $x = 4/3$, and $\mu_{23}(J(\QQ_{23}))$ is the group of order~$8$ generated by the
images of~$T_1$, $T_2$ and of the points of~$C(\QQ_{23})$ with $x = 0$.
\end{lemma}

\begin{proof}
\leanok
\uses{lem:ds-padic-squares}
As for Lemma~\ref{lem:ds-e1-local}: the orders are $\#J(\QQ_v)[2] = 2^4$ and~$2^3$
(Lemma~\ref{lem:ds-loc}\,(1) and Theorem~\ref{thm:ds-two-torsion}), as~$f$ has five, resp.~four,
irreducible factors over~$\QQ_v$ (Lemma~\ref{lem:ds-e2-padic}\,(2)). The points exist:
$f(-1/2) = 385/32$ and $f(4/3) = 11440/243$ are units of~$\ZZ_{59}$ with residues $6^2$ and~$4^2$,
and $f(0) = 12 \equiv 9^2 \bmod 23$. Their classes, and those of~$T_1$ and~$T_2$
(Corollary~\ref{cor:ds-torsion-values}), are compared through their components at the roots
of~$f$ in~$\QQ_v$, each mapped by Lemma~\ref{lem:ds-padic-squares}\,(2) to its valuation
modulo~$2$ and the class of its residue, computed by Lemma~\ref{lem:ds-e2-padic}\,(3). At~$59$
the four classes are independent. At~$23$ the classes of~$T_1$ and~$T_2$ are independent, and
the class~$\delta$ of the point with $x = 0$ has trivial components at the roots $-4$, $-3$
and~$r$, as $4$, $3$ and $-r \equiv 3$ are squares modulo~$23$; but its component $-\theta_h$
in~$E$ is not a square (Lemma~\ref{lem:ds-e2-padic}\,(4)), so the three classes are independent.

\emph{Formalization.} As in~E1, with three general lemmas:
\texttt{HypJac.LocalStructure.range\_descentMap\_eq\_of\_le\_of\_card\_ge} (for odd~$v$, a subgroup
of the local image of order at least $\#J(\QQ_v)[2]$ is all of it),
\texttt{HypJac.mk\_mem\_range\_descentMap\_of\_isSquare} (the class of $x_0 - \theta$ lies in the
local image if $f(x_0)$ is a square) and \texttt{Padic.isSquare\_iff\_modPowTwoCoords\_eq\_zero};
the lower bounds come from \texttt{Subgroup.card\_closure\_range\_eq\_two\_pow\_of\_monoidHom}. The
components at the roots are evaluations $L_v \to \QQ_v$ (\texttt{AdjoinRoot.liftAlgHom}), with the
coordinates \texttt{DescentExample2.rootBits}, and, for~$E$, the map
\texttt{AdjoinRoot.algHomOfDvd} to $\QQ_{23}[x]/(h)$; over~$\QQ_{23}$ no isomorphism of
$L_v^\times/(L_v^\times)^2$ with a product is needed. The number of factors over~$\QQ_{23}$ comes
from \texttt{Polynomial.Factors.card\_eq\_of\_associated\_prod}, over~$\QQ_{59}$ from
\texttt{Polynomial.Factors.card\_eq\_of\_injective} as in~E1. Only components of odd valuation
need the roots modulo~$p^2$: $-4 - r_1$ of~$\mu_{59}(T_1)$ and $-3 - r$ of~$\mu_{23}(T_2)$ here,
and the generators of the primes of Lemma~\ref{lem:ds-e2-field}\,(2) at their own roots in
Proposition~\ref{prop:ds-e2-sel}.
\end{proof}

\begin{proposition}[The relaxed Selmer group of~E2]
\label{prop:ds-e2-sel}
\lean{ChabautyColeman.DescentExample2.card_relaxedSelmerGroup_le}
\leanok
\uses{def:ds-relaxed, lem:ds-e2-selmer, lem:ds-e2-local}
Assume~(Loc) for~$\QQ_{23}$ and~$\QQ_{59}$. Then $\#S(\Sigma, V) \le 8$.
\end{proposition}

\begin{proof}
\leanok
\uses{thm:ds-naturality, lem:ds-e2-padic, lem:ds-padic-squares}
Each root~$\rho$ of~$f$ in~$\QQ_v$, for $v = 23, 59$, gives a homomorphism
$L^\times/(L^\times)^2 \to \QQ_v^\times/(\QQ_v^\times)^2 \to \FF_2^2$: the component at~$\rho$, then
its valuation modulo~$2$ and whether its residue is a nonsquare. Together, the roots $-4$, $-3$,
$r_1$, $r_2$, $r_3$ of~$f$ in~$\QQ_{59}$ and $-4$, $-3$, $r$ in~$\QQ_{23}$ give
\[ \beta = (\beta_{59}, \beta_{23}) \colon L^\times/(L^\times)^2 \to \FF_2^{10} \times
   \FF_2^6\,, \]
with~$\beta_v$ factoring through~$\operatorname{res}_v$ (Theorem~\ref{thm:ds-naturality}). By
Lemma~\ref{lem:ds-e2-local}, the images of the local images $\mu_v(J(\QQ_v))$ under~$\beta_{59}$
and~$\beta_{23}$ have dimensions~$4$ and~$2$, as the class~$\delta$ of the point with $x = 0$ lies
in the kernel of~$\beta_{23}$. So there are ten linearly independent linear forms on
$\FF_2^{10} \times \FF_2^6$ vanishing on both, which give a homomorphism
$\Theta \colon L^\times/(L^\times)^2 \to \FF_2^{10}$ with $S \subseteq \ker\Theta$. The forms can
be chosen so that~$\Theta$ maps the ten generators $(-1, 1, 1)$, $(23, 1, 1)$, $(1, 23, 1)$,
$(1, 59, 1)$ and $(1, 1, \varphi)$, $\varphi = t$, $t + 3$, $2t - 3$, $t + 4$, $5t + 6$,
$t^2 + 4t - 3$, of~$A(\Sigma, 2)$ (Lemma~\ref{lem:ds-e2-selmer}) to the standard basis
of~$\FF_2^{10}$. So $\Theta(A(\Sigma, 2)) = \FF_2^{10}$, and
$\#S \le \#(A(\Sigma, 2) \cap \ker\Theta) = \#A(\Sigma, 2)/2^{10} \le 2^3$.

\emph{Formalization.} The forms are a $16 \times 10$ matrix over~$\ZZ/2$, computed with Magma and
checked in Lean by \texttt{decide}: it kills the values of~$\beta$ on the generators of the local
images and maps the ten generators to the standard basis. The map~$\beta_v$ is
\texttt{DescentExample2.globalBits}, with values in~$(\ZZ/2)^2$ per root
(\texttt{Padic.modPowTwoEquivZModTwo}; the second coordinate is the quadratic character
of~$\ZZ/p$, Mathlib's \texttt{quadraticChar}, through
\texttt{FiniteField.modPowTwoEquivZModTwo}). The count runs in the subgroup~$B$ generated by the
$13$ classes, which contains $A(\Sigma, 2)$ (\texttt{DescentExample2.selmerGroup\_le\_closure}), so
the classes need not be shown to lie in~$A(\Sigma, 2)$: $\#(B \cap \ker\Theta) \cdot
\#\Theta(B) = \#B \le 2^{13}$, by Mathlib's \texttt{Subgroup.relIndex\_mul\_card} and
\texttt{Subgroup.relIndex\_ker}, and \texttt{Subgroup.card\_closure\_range\_le\_two\_pow}. The norm
condition is not used, and the remaining three generators $(59, 1, 1)$, $(1, -1, 1)$,
$(1, 1, -1)$ need no values.
\end{proof}

\begin{lemma}[$P$ has infinite order]
\label{lem:ds-e2-points}
\lean{ChabautyColeman.DescentExample2.not_isOfFinOrder_P,
  ChabautyColeman.DescentExample2.card_closure_descentMap_eq}
\leanok
\uses{thm:ds-formula, cor:ds-torsion-values, thm:ds-naturality, thm:ds-two-torsion,
  lem:ds-e2-padic}
The classes $\mu(T_1)$, $\mu(T_2)$ and $\mu(P)$ generate a subgroup of order~$8$
of~$L^\times/(L^\times)^2$, and~$P$ has infinite order.
\end{lemma}

\begin{proof}
\leanok
\uses{lem:ds-padic-squares}
The values are (Corollary~\ref{cor:ds-torsion-values} and Theorem~\ref{thm:ds-formula})
\[ \mu(T_1) = (59, -1, -4 - t)\,, \qquad \mu(T_2) = (1, -23, -3 - t)\,, \qquad
   \mu(P) = (25/9, 16/9, -11/9 - t)\,. \]
Their images under the map~$\beta$ of the proof of Proposition~\ref{prop:ds-e2-sel} are
independent: only $\mu(T_1)$ has odd valuation at the root~$-4$ in~$\QQ_{59}$; $\mu(P)$ has trivial
components in~$\QQ_{23}$, unlike~$\mu(T_2)$; and the component $-11/9 - r_2$ of~$\mu(P)$ is a unit
of~$\ZZ_{59}$ that is not a square modulo~$59$. If~$P$ had finite order,
$B = J(\QQ)[2] + \langle P\rangle$ would be a finite group with $B[2] = J(\QQ)[2]$, of order~$4$
(Theorem~\ref{thm:ds-two-torsion}, as~$f$ has three irreducible factors). As~$\mu$ kills~$2B$ and
$\#(B/2B) = \#B[2]$ for a finite abelian group, $\#\mu(B) \le 4$, whereas~$\mu(B)$ contains the
group of order~$8$ above.

\emph{Formalization.} Neither~(MW) nor~(Loc) is used. The counting step is
\texttt{Subgroup.card\_map\_le\_card\_inf\_ker\_powMonoidHom}: for a finite subgroup~$B$ of a
commutative group~$G$ and a homomorphism~$\varphi$ to a group of exponent~$n$,
$\#\varphi(B) \le \#(B \cap G[n])$, here with $n = 2$; a subgroup generated by finitely many
elements of finite order is finite (\texttt{Subgroup.finite\_closure\_of\_isOfFinOrder}, from
Mathlib's \texttt{CommGroup.finite\_of\_fg\_isMulTorsion}). The Lean takes for~$B$ the subgroup
generated by~$T_1$, $T_2$ and~$P$ and bounds $\#B[2]$ by $\#J(\QQ)[2] = 4$, so it does not need
$J(\QQ)[2] = \langle T_1, T_2\rangle$.
\end{proof}

\begin{theorem}[E2]
\label{thm:ds-e2}
\lean{ChabautyColeman.DescentExample2.finrank_hypJac,
  ChabautyColeman.DescentExample2.relaxedSelmerGroup_eq}
\leanok
\uses{def:ds-cited, thm:ds-rank-bound, prop:ds-e2-sel, lem:ds-e2-points}
Assume~(MW) and~(Loc) for~$\QQ_{23}$ and~$\QQ_{59}$. Then $\operatorname{rk} J(\QQ) = 1$, and
$S(\Sigma, V) = \mu(J(\QQ)[2] + \langle P\rangle)$, of order~$8$.
\end{theorem}

\begin{proof}
\leanok
\uses{thm:ds-two-torsion, def:ds-relaxed}
By Theorem~\ref{thm:ds-rank-bound} and Proposition~\ref{prop:ds-e2-sel},
$2^{\operatorname{rk} J(\QQ)} \cdot 4 \le 8$, as $\#J(\QQ)[2] = 2^{3 - 1}$
(Theorem~\ref{thm:ds-two-torsion}). By Lemma~\ref{lem:ds-e2-points}, $P$ has infinite order, so
$\operatorname{rk} J(\QQ) \ge 1$, as $J(\QQ)$ is finitely generated. Finally
$\mu(J(\QQ)[2] + \langle P\rangle) \subseteq S$ (Definition~\ref{def:ds-relaxed}) has order~$8$ by
Lemma~\ref{lem:ds-e2-points}.

\emph{Formalization.} The rank is \texttt{Module.finrank} over~$\ZZ$ of \texttt{Additive (HypJac
$\QQ$ f)}; the hypotheses are the instance that \texttt{HypJac $\QQ$ f} is finitely generated (which
\texttt{HypJac.MordellWeil} provides) and \texttt{HypJac.LocalStructure} for~$\QQ_{23}$
and~$\QQ_{59}$; the statement on $S(\Sigma, V)$, \texttt{relaxedSelmerGroup\_eq}, needs~(Loc) only.
The instances of Theorem~\ref{thm:ds-rank-bound} on the class groups and unit groups of the field
factors are found for $K = \QQ$ in general, as the instances
\texttt{NumberField.instFiniteClassGroupIntegralClosure} and
\texttt{NumberField.instIsMulFGUnitsIntegralClosure}; $\operatorname{rk} J(\QQ) \ge 1$ comes through
Mathlib's \texttt{Module.finrank\_eq\_zero\_iff\_isTorsion}. The order~$8$ of $S(\Sigma, V)$ is not
part of the Lean statement; it is \texttt{DescentExample2.card\_closure\_descentMap\_eq} with the
proof.
\end{proof}
